\documentclass[../main.tex]{subfiles}
\begin{document}
%==========================================TIÊU ĐỀ TẬP========================================
\begin{table}[h]
  \begin{adjustwidth}{-1cm}{}
    \begin{tabular}{>{\centering\arraybackslash}p{6cm}|>{\raggedright\arraybackslash}p{12.5cm}}
      \multirow{5}{6cm}
      % Cột bên trái: ÉPISODE và số tập
      {\\[-40pt]
        \flaregothic\fontsize{20pt}{20pt}\selectfont \centering
        {\contour{errortwo}{\textcolor{black}{ÉPISODE}}}\color{black}\\
        \burbank\fontsize{72pt}{72pt}\selectfont
        \includegraphics[height=3cm]{../../../../Image/Book?/number/num8.png}
        \\[18pt]
        % \flaregothic\fontsize{14pt}{14pt}\selectfont
        % \color{epattr}BIENVENUE, \uline{\color{Banpire} Mel} !
        % \flaregothic\fontsize{18pt}{18pt}\selectfont
        % \color{epattr}ÉPISODE PERDU
        \color{black}
      }
       &
      % Dòng thứ nhất: tên tập tiếng Nhật
      {\fotskip\fontsize{18pt}{18pt}\selectfont また、\ruby{町}{まち}へ
      }                                                                                            \\[6pt]
      % Dòng thứ hai: tên tập tiếng Anh
       & {\cafeteria\fontsize{18pt}{18pt}\selectfont The Journey Back}                                                                                     \\[4pt]
      % Dòng thứ ba: tên tập tiếng Việt
       & {\vnmsans\fontsize{18pt}{18pt}\selectfont Lời ước nguyện dang dở}                                                                                 \\[8pt]
      % Dòng thứ tư: Nhân vật xuất hiện
       & {\caviar{\fontsize{14pt}{14pt}\selectfont Track used: The Apocalypse (Last Reunion - Peter Roe) (from Dancing Line: Community Edition/Max Line)}}
      \\[6pt]
      % Dòng cuối cùng: Thời gian diễn ra của tập (thường sẽ là ngày phát hành)
       & {\continuum\fontsize{16pt}{16pt}\selectfont Ngày phát hành: 08/04/2022}                                                                           \\[8pt]
    \end{tabular}
  \end{adjustwidth}
\end{table}
%===========================================NỘI DUNG==========================================
\color{black}

\pov{Haragen Reiko}

``Hộc hộc\dots''

Tôi thở dốc, từng tiếng khò khè như nghẹn lại vì bụi. Cái không khí mặn mặn của đất đá xộc vào mũi, lẫn mùi gỗ cháy và mùi kim loại, như thể cả ngôi trường vừa trải qua một cơn sốt kinh hoàng và giờ chỉ còn lại hơi thở yếu ớt của các bức tường. Tôi hắng giọng, cố gắng nuốt cơn ho cho khỏi làm vỡ đi cái khoảng tĩnh lặng đau nhói đang bao quanh.

``Mọi người\dots\ Có ai không!? Shino-san! Kyuen-san! Saikyou-san! \textbf{Các cậu ở đâu!?}'' tiếng tôi vang lên hoang vắng giữa những hàng ghế vỡ vụn, bị dội lại bởi mảng tường đổ và khung cửa bị sập. Mỗi âm thanh như rơi vào lòng đất, bị nuốt chửng bởi sự yên lặng đột ngột.

``\textbf{Làm ơn trả lời mình đi!}'' tôi gào lên, giọng khàn đặc. ``\textbf{Ai đó\dots\ làm ơn!}''

Tôi chạy, chân đạp lên mảnh gạch vỡ, văng ra từng tiếng lộc cộc; đầu óc quay cuồng vì bụi bay vào mắt, tôi phải dùng tay quơ đi để còn thấy được đường.

``Khụ khụ\dots\ Yae-san? Sakura-san? Shiina-san! \textbf{Các cậu đâu rồi!?}''

Lớp học 1-B, nơi mà chỉ vài giờ trước còn tiếng nói cười của các bạn, giờ đây la liệt những mảnh gỗ vỡ, các giá bảng ngang nghiêng, những mảng trần gỗ lửng lơ, vùi lấp trong đống đất đá. Mọi thứ im lặng một cách khủng khiếp.

``\textbf{Xin các cậu\dots\ đừng làm mình sợ!}'' tôi vừa chạy vừa gọi, giọng run rẩy. ``\textbf{Honoka-san! Nanase-san! Các cậu ở đâu!?}''

Tim tôi đập thình thịch khi len lỏi qua từng đống đổ nát. Mỗi bước chân như dẫm lên một nỗi sợ hãi mới. Tôi cố không nghĩ đến điều tồi tệ nhất, nhưng những hình ảnh kinh hoàng cứ hiện lên trong đầu.

``Làm ơn\dots\ làm ơn đấy\dots'' tôi thì thầm như một lời cầu nguyện.

Bất chợt, tôi dừng khựng lại. Dưới đống đổ nát gần cửa sổ, tôi thấy hai bóng người ôm chặt lấy nhau. Tim tôi như ngừng đập. Shiina-san và Honoka-san nằm đó, vẫn trong tư thế ôm nhau như lúc cuối cùng, và bên cạnh họ là Yae, tay vẫn nắm chặt lấy áo của cô bạn tóc vàng.

``Không\dots\ không\dots'' tôi lảo đảo tiến lại gần. ``\textbf{Shiina-san! Honoka-san! Yae-san!}''

Tôi quỳ xuống bên cạnh họ, tay run rẩy chạm vào má Shiina. Lạnh ngắt. Mắt cô vẫn mở to, ánh nhìn đông cứng trong nỗi kinh hoàng. Honoka-san nằm gục trong vòng tay cô, những giọt nước mắt còn đọng trên mi. Yae nằm bên cạnh, khuôn mặt hiền hậu giờ tái nhợt như sáp.

``\textbf{Tỉnh dậy đi!}'' tôi lay vai họ, nước mắt trào ra. ``Làm ơn\dots\ Hai người tỉnh dậy đi\dots''

Nhưng không có phản ứng nào. Chỉ có sự im lặng chết chóc đáp lại. Tôi gục xuống, từng tiếng nấc thốt ra trong cổ họng.

``Tại sao\dots\ tại sao lại thế này\dots''

Cô ấy quặp mình lại, một tay ôm lấy đầu như thể cố níu giữ những thứ đang nghiền nát trong đầu. Gương mặt cô méo xệch vì đau đớn. Tôi gập người xuống, kêu tên cô, đưa tay chạm lên trán cô, thật lạnh. Một vết lún sâu ở thái dương, hoặc có thể một mảng xi măng đã đè lên, khiến mặt cô méo mó. Khi tôi lật cô thật nhẹ, con ngươi vẫn mở nhưng đã vô hồn.

Từ trong ngực tôi, một tiếng thét nho nhỏ bật lên, rồi nghẹn thành nước.

``Nanase-san\dots'' tôi quỳ xuống bên cạnh cô, nước mắt rơi xuống khuôn mặt méo mó của cô. ``Xin lỗi\dots\ xin lỗi vì mình đến muộn\dots''

Tôi vẫn chưa chấp nhận nổi, không thể chấp nhận được sự thật này. Rằng những người bạn học cùng lớp mà tôi mới làm quen được ít lâu, những con người còn đang cười nói, đùa giỡn cách đây vài tiếng đồng hồ - Yae-san với nụ cười hiền hậu, Shiina-san luôn mạnh mẽ, Honoka-san hay đùa vui, và Nanase-san thông minh - giờ đây đã vĩnh viễn rời xa tôi, không bao giờ còn có thể mở mắt ra nữa.

``\textbf{Tại sao!?}'' tôi thét lên, giọng vỡ òa. ``\textbf{Tại sao lại là các cậu!? Tại sao!?}''

Tôi nhớ lại buổi chiều hôm đó, khi Honoka-san rủ cả nhóm thử thách vào phòng Mỹ thuật - nơi đồn đại có ma. ``Chúng ta sẽ tìm xem liệu lời đồn đại của mọi người có đúng hay không nha\~{}!'', cô ấy nói vậy. Tôi đã từ chối, chỉ đứng bên ngoài trường chờ kết quả. Nhưng các bạn khác hào hứng nhận lời, thậm chí hai anh em ấy dù có chút do dự nhưng cũng đồng ý.

Ai mà ngờ được, chỉ vài tiếng sau, trận sạt lở đất kinh hoàng kia đã vùi lấp tất cả.

Cảm giác tan nát không đến một lúc, mà dần dần lan ra, như lệ tràn qua từng ngón tay. Tôi ngồi bệt xuống, ôm đầu gối, nước mắt trào ra không ngừng. ``Giá như mình ở lại\dots\ giá như mình ngăn họ lại\dots'', lời hối tiếc thất thanh của tôi bị nuốt chửng bởi tiếng gió thổi qua những lỗ thủng trên mái.

``Tại sao\dots\ tại sao lại như vậy?'' tôi lặp lại câu hỏi mà chính tôi cũng không thể trả lời.

Tôi cố gắng cứng rắn, đứng dậy, lau mặt, nuốt cơn ói lên cổ họng vì mùi bụi. Bốn thi thể trước mặt tôi, bốn cái tên mà tôi vừa mới gọi vang vọng, giờ yên lặng như những bức tượng. Tôi lần lượt nắm tay từng người, đặt tai lên tim họ như lúc kiểm tra nhịp cho một đứa trẻ: không thấy gì cả. Mỗi cái chạm làm tôi run hơn, xúc giác biến thành một nỗi đau lạnh lùng.

``\textit{Còn những người khác\dots}'' tôi lẩm bẩm, cố gắng đứng vững. ``\textit{Mình phải tìm những người khác\dots}''

Cơn tuyệt vọng tràn ngập, nhưng tôi không có quyền gục xuống ở đây. Tôi buộc mình đứng dậy, chân tiếp bước chạy đi, miệng gọi tiếp các bạn học khác: ``Shino-san! Sakura-san! Kyuen-san! Saikyou-san!'' Từng tiếng gọi như rơi xuống một cái hố không điểm dừng.

``\textit{Làm ơn\dots\ hãy còn sống\dots}'' tôi thì thầm trong nước mắt. ``\textit{Xin các cậu\dots\ hãy còn sống\dots}''

Tôi chạy tiếp, bước qua những hàng ghế méo mó, qua bàn vẽ nghiêng ngả, lên những bậc cầu thang gãy. Mỗi tiếng gọi như một lần lột tả nỗi sợ: họ có thể ở ngay phía sau tấm ván kia, dưới một tấm thảm vải, ở trong phòng y tế, bất cứ đâu. Tim tôi như một chiếc búa muốn bật khỏi lồng ngực, nhen nhóm một hy vọng nhỏ nhoi cho những bạn học còn lại.

Tôi không biết mình còn có thể bền được bao lâu. Mỗi bước chân là một quyết định giữa ``không thể ở lại đây'' và ``phải tìm họ bằng mọi giá''. Tôi lẳng lặng hít một hơi sâu, cố gạt nước mắt ra để còn thấy đường.

Trong khoảnh khắc yên ắng đến nghẹt thở ấy, một âm thanh lạ lọt vào tai tôi. Tôi khựng lại, tim đập thót lên một nhịp. Có phải tôi vừa nghe thấy gì đó? Hay chỉ là tiếng gió thổi qua những khe nứt? Tôi nín thở, lắng nghe, cố phân biệt giữa tiếng đổ nát rơi rụng và những âm thanh khác.

``Ư\dots\ đau quá\dots'' một tiếng rên yếu ớt, đứt quãng vang lên, nhỏ đến mức tưởng như có thể tan biến vào không khí, nhưng nó thật sự tồn tại.

Tôi giật bắn người, toàn thân run lên dữ dội. Một cảm giác vừa hoảng loạn vừa mừng rỡ tràn ngập trong lồng ngực - hoảng loạn vì đó chắc chắn là tiếng người, mừng rỡ vì điều đó đồng nghĩa với việc còn người sống sót.

``\textit{Có người!}'' tôi thì thầm, giọng run rẩy. ``\textit{Thật sự có người!}'' tôi quay đầu lia lịa, cố định vị hướng phát ra âm thanh.

Ở phía hành lang, cách bốn cái xác không xa, có một mảng gạch đá vẫn còn âm ỉ khói bụi. Tôi tiến lại gần, tim đập liên hồi, và rồi\dots\ tôi thấy nó. Một cánh tay chìa ra từ dưới lớp tường gạch, ngón tay cử động yếu ớt. Trên cổ tay ấy, tôi nhìn rõ một vòng tay đỏ với ba chiếc chuông bạc nhỏ lấp lánh giữa bụi mờ. Tôi chết lặng. Vòng tay này chính là thứ Saikyou-san đã đưa cho người anh của cô, và sau đó cậu ấy giữ lấy. Tôi nhớ rõ, không thể nhầm được.

Vậy thì người này là\dots

``\textbf{Kyuen-san!}'' tôi hét lên, quỳ sụp xuống, không để tâm gì đến những viên gạch sần sùi cào rách da tay. Tôi vừa gọi tên vừa điên cuồng bới từng mảng gạch vụn. Lực tay của tôi không lớn, ngón tay run rẩy, vụng về và vội vã, nhưng tôi không quan tâm nữa. Mỗi lần nhấc được một mảnh ra, tôi lại càng sốt ruột.

``Ráng lên, Kyuen-san! Đừng bỏ mình ở đây một mình\dots\ xin cậu, đừng giống họ!''

Cuối cùng, dưới lớp gạch xám lẫn những thanh gỗ dập nát, một tấm lưng hiện ra. Lưng của một chàng trai, đầy bụi đất, chiếc áo đen rách bươm, máu rịn đỏ từ thái dương chảy xuống. Tôi run rẩy đặt tay lên, gọi: ``Kyuen-san! Nghe tôi không?''

Cậu ta khẽ động đậy. Tiếng rên bật ra, mơ hồ: ``Ư\dots\ Reiko\dots\ san\dots?'' Môi cậu run rẩy, giọng khàn đặc. Tôi muốn òa khóc ngay lúc đó. Kyuen-san vẫn còn sống. Khác với Yae, Shiina, Honoka, Nanase đã bất động, cậu vẫn còn ở đây, ngay trước mắt tôi.

Nhưng rồi tôi sững lại. Bởi vì tôi nhận ra cậu không nằm một mình. Hai tay cậu, yếu ớt mà kiên quyết, đang ôm ghì một thân người khác bên dưới.

``Cậu\dots\ cậu ôm ai vậy\dots?'' tôi lắp bắp, tim đập dồn.

Kyuen gượng gạo lật người sang bên. Dưới tấm thân bê bết bụi của cậu, tôi thấy một cô gái có nét mặt tương đồng, mái tóc dài vướng đất đá, trên tay và gương mặt đầy vết xước, nhưng đôi mắt xanh dương vẫn mở hé. Saikyou-san. Cô khụt khịt ho, từng cơn ho dồn dập vì khói bụi, nhưng vẫn còn thở.

``Reiko-san\dots\ Là\dots\ cậu à\dots?'' cô ấy thều thào, giọng run và khàn, nhưng là minh chứng cho việc cô còn sống.

Cả người tôi mềm nhũn, đầu gối khuỵu xuống. Tôi gần như gào lên trong sự nhẹ nhõm không thể diễn tả: ``Tốt quá rồi\dots\ \textbf{Trời ơi, tốt quá rồi!}''

Tôi không kìm được, vội ôm chầm lấy cả hai, mặc cho bụi đất bám dày trên quần áo và mùi máu tanh bám trên tay áo Kyuen-san. Cảm giác ôm được hai thân người còn hơi ấm khiến nước mắt tôi lại rơi như mưa, lần này không chỉ vì đau đớn, mà còn là sự nhẹ nhõm khi cuối cùng tôi vẫn còn người sống.

Kyuen thở hổn hển, cố mỉm cười gượng gạo. ``May mà\dots\ cậu tìm được bọn tôi\dots\ Nếu không thì\dots''

Saikyou cũng run rẩy, tay nhỏ bé bấu lấy tay tôi. ``Reiko-san\dots\ tôi tưởng\dots\ tôi tưởng sẽ không bao giờ\dots''

Tôi ghì chặt họ, thốt lên: ``Không sao rồi\dots\ hai người còn sống, thế là đủ rồi\dots\''

Nhưng rồi, sự thực cay đắng vẫn treo lơ lửng trên đầu tôi. Khi niềm vui ban đầu dần lắng xuống, cô em gái của Kyuen-san thở gấp, hỏi trong vô thức: ``Những người khác\dots\ cậu đã gặp họ chưa?''

Tim tôi thắt lại. Tôi im lặng, không dám trả lời. Câu hỏi ấy như một nhát dao khoét sâu vào vết thương chưa khép. Cậu con trai nhìn thẳng vào mắt tôi, đôi mắt mệt mỏi nhưng kiên định. ``Reiko-san\dots\ đừng im lặng như thế. Nói đi. Họ\dots\ họ đâu rồi?''

Nước mắt tôi lại trào ra, nghẹn ngào. ``Họ\dots\ không còn nữa.'' Tôi cắn môi, cố gắng dứt khoát, chỉ về phía hành lang nơi bốn cái xác nằm lại. ``Yae-san, Shiina-san, Honoka-san, và cả\dots\ Nanase-san\dots\ tất cả\dots\ đã đi rồi.''

Không gian như đông cứng.

Kyuen-san chết lặng, mắt trợn trừng, hơi thở hụt hẫng, như thể cậu đã đoán định được. Cậu nghiến chặt răng, nắm tay run lên hiện rõ sự bất lực, máu từ vết thương lại rịn ra nhưng cậu không còn cảm thấy.

Saikyou-san thì run bắn, hai tay ôm miệng, giọt lệ trào ra không kịp ngăn. ``Không thể nào\dots\'' cô lắc đầu liên tục, như muốn phủ nhận thực tại.

Tôi thấy đôi mắt họ thoáng nhòa đi, rồi cả hai đồng thời nhìn về phía bốn thân người nằm bất động. Một ý nghĩ chợt lóe lên trong đầu tôi.

``Hai người\dots\ đến từ tương lai, phải không?'' tôi hỏi, giọng có chút quở trách. ``Nếu đã biết trước chuyện này sẽ xảy ra, tại sao không nói gì cả!?''

Cả hai anh em im lặng một lúc. Cuối cùng, người em gái khẽ lên tiếng, giọng run run: ``Mọi chuyện\dots\ diễn ra quá nhanh\dots''

``Không ai trong chúng tôi kịp phản ứng cả,'' Kyuen tiếp lời, giọng đau đớn. ``Dù có biết trước\dots\ nhưng khi thảm họa thực sự ập đến\dots''

Cậu không nói hết câu, nhưng tôi hiểu. Đôi khi, ngay cả khi biết trước tương lai, con người vẫn bất lực trước những điều không thể thay đổi.

Cả ba chúng tôi đứng lặng trước cảnh tượng đau thương. Tôi nắm chặt tay Saikyou-san đến mức các khớp ngón tay trắng bệch, như thể sợ rằng nếu buông ra, cô ấy cũng sẽ tan biến như những người bạn kia. Bên cạnh, Kyuen-san gượng người ngồi dậy, đôi mắt đỏ hoe nhìn chăm chăm vào bốn thân thể bất động, như không dám tin vào thực tại trước mắt.

Không ai nói được lời nào. Chỉ có tiếng thở nặng nề và những giọt nước mắt lặng lẽ rơi. Trước mắt chúng tôi, bốn con người từng tràn đầy sức sống - Yae-san với nụ cười ấm áp và lý trí người con gái, Shiina-san luôn mạnh mẽ kiên cường, Honoka-san tinh nghịch vui vẻ, và Nanase-san kiêu kỳ nhưng thông minh điềm đạm - giờ đây chỉ còn là những hình hài cứng đờ, lạnh lẽo. Những ước mơ, những lời hứa về tương lai, về những ngày học tập cùng nhau, tất cả đều tan vỡ dưới bàn tay tàn nhẫn của số phận.

Saikyou-san run rẩy quỳ xuống bên cạnh họ, bàn tay nhỏ bé của cô vươn ra chạm vào má của cô gái tóc vàng, rồi rụt lại như bị bỏng. ``Không\dots\ không phải thế này\dots'' cô thì thầm, giọng vỡ vụn. Kyuen-san đặt tay lên vai em gái, nhưng chính bờ vai cậu cũng đang run lên từng đợt. Tôi thấy rõ cậu đang cố kìm nén cơn đau đang xé nát lồng ngực.

Tôi quỳ xuống bên cạnh họ, cổ họng nghẹn đắng. Những kỷ niệm về buổi sáng hôm nay - tiếng cười đùa trong lớp học, những câu chuyện vui vẻ trong giờ giải lao, lời hẹn gặp nhau vào ngày mai - giờ như những mũi kim đâm vào tim. ``Xin lỗi\dots'' tôi thì thầm, nước mắt rơi xuống nền đất lạnh, ``xin lỗi vì đã không ở lại\dots\ xin lỗi vì đã không thể cứu các cậu\dots''

Saikyou-san bật khóc thành tiếng, những tiếng nấc vang vọng trong không gian đổ nát. Kyuen-san ôm chặt em gái vào lòng, cậu không khóc, nhưng tôi thấy rõ những giọt lệ cậu đang cố nuốt ngược vào trong, và đôi mắt cậu chất chứa một nỗi đau không gì có thể xoa dịu.

Chúng tôi không thể ở lại đây mãi. Tôi nắm tay Saikyou kéo dậy, rồi đỡ Kyuen đứng lên, cậu khập khiễng nhưng vẫn gượng. Cả ba quay lưng lại, từng bước nặng trĩu rời khỏi nơi ấy.

``Còn Sakura-san\dots\ còn Shino-san\dots'' tôi lẩm bẩm, giọng vỡ vụn. Hai người bên cạnh khẽ gật, không ai nói thêm gì.

Chúng tôi bước đi, hòa trong bóng tối đặc quánh, chỉ còn một tia hy vọng mỏng manh dẫn đường: rằng hai người kia\dots\ vẫn còn ở đâu đó, đang chờ được tìm thấy.

\ct

Bụi đất vẫn còn lơ lửng trong không khí, mỗi bước đi của chúng tôi như dẫm lên những tàn dư nóng hổi của một cơn ác mộng vừa trút xuống. Tôi siết chặt bàn tay Saikyou-san, còn cánh tay kia vòng qua người Kyuen-san, cố giữ cậu không ngã. Vết thương trên đầu cậu ấy vẫn rỉ máu, sắc mặt cậu trắng bệch đến mức tôi sợ chỉ cần một cơn gió mạnh cũng đủ khiến cậu gục xuống ngay tức khắc.

``Reiko-san, cậu đỡ anh ấy kỹ một chút,'' Saikyou-san thì thầm, giọng căng thẳng. Cô vòng người sang, ôm lấy cánh tay còn lại của người anh trai, nửa dìu nửa kéo.

Người anh cố gượng cười, đôi môi nứt nẻ, run rẩy nhưng vẫn cứng cỏi: ``Anh\dots\ vẫn đi được. Đừng lo cho anh làm gì.''

Cô em gái của cậu cau mày, ngước lên nhìn anh trai, giọng lạc đi: ``Đầu anh nứt mẻ cả trán thế kia mà bảo là không sao thật ư!?''

``Thật mà,'' cậu vẫn khăng khăng, nhưng chỉ sau vài bước, chân cậu lảo đảo, suýt nữa ngã dúi dụi. Tôi và Saikyou-san phải gấp rút giữ chặt, nếu không cậu đã nằm sõng soài trên đống gạch vụn.

``Anh còn đi kiểu này thì tụi em chẳng thể nào tìm được Sakura-san với Shino-san đâu, Nii-chan!'' cô em nén nước mắt, cắn môi đến bật máu, rồi dứt khoát: ``Nii-chan, leo lên lưng em đi.''

``Không cần\dots'' Kyuen khẽ lắc đầu, giọng yếu ớt.

``\textbf{Anh mà ngã xuống thêm lần nữa thì em thề em đánh anh liền đó!}'' cô gào lớn, đôi mắt đỏ hoe. Không khí như đông đặc lại một thoáng. Cũng dễ hiểu vì sao cô ấy lại phản ứng như vậy, khi Kyuen-san đã dùng chính thân mình để che chắn cho cô em gái, và có lẽ trông cậu đã ở bên bờ vực sụp đổ.

Kyuen nhìn em gái, mệt mỏi thở dài. ``Anh không muốn làm gánh nặng--''

``\textbf{Gánh nặng cái đầu anh!}'' Saikyou liền ngắt lời cậu, nước mắt trào ra. ``Anh mà mất đi nữa thì em sống sao được\dots''

Nhìn những giọt lệ trào ra từ một người suýt mất đi người mà cô thân yêu nhất, tôi không khỏi áy náy. ``Kyuen-san, cậu nên nghe lời em gái đi,'' tôi khẽ nói. ``Nhìn cậu lúc này, mình thấy lo lắm. Vết thương trên đầu vẫn chảy máu kìa. Với lại\dots''

Tôi nhìn gương mặt của Saikyou-san, hiện lên một sự cương quyết như một vị phụ huynh muốn người con phải nghe lời. ``\dots mình không nghĩ em gái cậu sẽ không chấp nhận câu trả lời `không' đâu.''

Người anh khựng lại, đôi môi run lên. Cuối cùng, cậu gục đầu xuống, chấp nhận để em gái đỡ mình trèo lên lưng. Cô em gái cắn răng, ghì chặt đôi chân gầy của anh trai, chật vật đứng thẳng dậy. Tôi nhìn cảnh đó, tim thắt lại. Hình ảnh ấy vừa bi thương vừa ấm áp đến mức tôi không dám quay mặt đi.

``\textit{Xin lỗi nhé, Kyuen-san, mình biết là cậu muốn tỏ ra mạnh mẽ trước em gái, nhưng\dots\ giờ không phải lúc đâu.}''

Chúng tôi tiếp tục bước. Lớp bụi dưới chân xào xạc, mỗi tiếng động vang lên trong không gian hoang tàn như một lời nhắc nhở về cái chết đang phủ quanh.

Sau một hồi im lặng, chính Saikyou-san mở lời, giọng khàn khàn vì sức nặng trên lưng: ``Nii-chan\dots\ Ngôi trường này\dots\ chỉ trong chớp mắt đã biến mất\dots\ Em còn tưởng\dots\ em mất anh luôn rồi.''

Kyuen-san trên lưng em khẽ cựa, cố gắng nói trong tiếng thở nặng nhọc: ``Anh biết\dots\ suýt nữa thì\dots\ Nhưng mà\dots\ còn em nữa. Anh phải che cho em chứ\dots''

Cả ba chúng tôi im lặng suốt một quãng dài. Chỉ còn tiếng bước chân kéo lê, tiếng ho sặc của người em gái xen lẫn tiếng thở hổn hển của anh trai. Nỗi đau của sự mất mát cứ thế dồn lại, đè nặng lên vai cả ba.

Rồi bất chợt, tôi nhìn thấy ở chỗ gạch vụn lẫn bóng tối khuất phía hành lang, có một bóng người nhỏ bé nằm vắt ngang trên đống đất đá. Tóc nâu lòa xòa phủ lên gương mặt, mái tóc ấy được cài một chiếc kẹp màu đỏ lệch sang bên trái.

Bóng dáng kia\dots\ không thể nhầm lẫn vào đâu được. Người bạn thân nhất của tôi.

``\textbf{Shino-san!}'' tôi hét lên, lao về phía trước. Kyuen-san trên lưng ngẩng phắt đầu, còn Saikyou-san thì lập tức tăng tốc, ba chúng tôi cùng hướng về phía bóng dáng ấy.

Và ở đó, chúng tôi nhìn thấy Shino-san nằm bất động, đôi vai nhỏ bé im lìm trong bụi đất.

\img{e1-7l}

Tôi lao vội về phía bóng dáng quen thuộc, đầu óc trống rỗng chỉ còn một niềm tin mong manh: đó là Shino-san, chắc chắn là Shino-san rồi.

Tôi quỳ sụp xuống bên cạnh, run rẩy đặt tay lên vai cô bạn và lay mạnh: ``Shino! Tỉnh lại đi\dots!''

Đôi mắt cô vẫn nhắm nghiền, gương mặt phủ đầy bụi đất, hơi thở thì\dots\ dường như chẳng còn chút dấu hiệu nào. Tôi như ngộp thở, bàn tay càng lúc càng mạnh hơn, gần như tuyệt vọng.

``Không phải cả cậu nữa\dots''

Ở phía sau, hai anh em chỉ đứng lặng, không tiến lại gần, vì họ biết rõ, nếu Shino-san thực sự đã rời bỏ, thì tôi sẽ là người đầu tiên cần được chuẩn bị tinh thần.

Nhưng tôi không muốn như vậy. Tôi không muốn mất thêm một người bạn nào nữa.

``Đừng\dots\ đừng giống như bốn người họ\dots\ đừng bỏ mình lại chứ\dots'' tôi thở rít, giọng nghẹn lại trong cổ.

Tôi lay người cô mạnh hơn, nhưng vẫn không có phản ứng. Không thể để mất thêm một người nữa, tôi vội vàng lật ngửa Shino ra, đặt tai lên ngực cô - nhịp tim yếu ớt, hơi thở mong manh. Không chần chừ, tôi bắt đầu ép tim và thổi ngạt cho cô, những động tác hô hấp nhân tạo mà tôi đã học được.

``Tỉnh lại đi, Shino-san!'' tôi vừa làm vừa thúc giục. ``Làm ơn\dots\ tỉnh lại đi!''

Và rồi, thật bất ngờ, đôi mí kia khẽ run lên. Một hơi thở nhẹ thoát ra, rất yếu nhưng đủ để tôi nhận ra sự sống chưa hoàn toàn tắt. Đôi mắt từ từ mở ra.

``Ưm\dots\ Tại sao\dots\ mình lại ở đây?'' một giọng nói nhỏ nhẹ vang lên.

\img{e1-8a}

Tôi bật khóc ngay tức thì, ôm chầm lấy cô bạn, siết thật chặt như thể nếu tôi buông ra thì cô sẽ biến mất mãi mãi. ``Shino-san! Mình\dots\ mình cứ tưởng cậu đã chết rồi chứ\dots''

Nhưng cơ thể trong vòng tay tôi hơi cứng lại. Người ấy không đáp lại vòng tay ôm của tôi, chỉ ngơ ngác nhìn xuống. ``Ơ\dots\ xin lỗi\dots\ Cậu\dots\ là ai vậy?''

Tôi chết lặng. Nước mắt còn chưa kịp ngừng, giờ lại đóng băng ngay trên má. ``Cậu\dots\ nói gì cơ? Mình là Reiko-san mà! Là Haragen Reiko đây!''

Người trước mặt tôi vẫn giữ gương mặt bối rối, rồi như chợt nhớ ra điều gì, cô ấy nói nhỏ: ``Haragen\dots? Tớ có biết hai người họ Haragen\dots\ nhưng chưa từng nghe ai tên Reiko cả\dots\'' Cô ngập ngừng một lúc rồi tiếp: ``Cậu\dots\ có nhầm không? Tớ thực sự chưa từng gặp cậu bao giờ\dots''

Câu nói ấy như một nhát dao lạnh ngắt xuyên vào tim tôi. Tôi lùi ra một chút, run rẩy nhìn thẳng vào đôi mắt cô. ``Shino-san\dots? Cậu đang nói gì vậy\dots? Đừng đùa với mình\dots''

Từ phía sau, tôi thấy Saikyou-san và Kyuen-san đã bắt đầu thì thầm với nhau. Họ nhìn kỹ hơn, và cả hai đồng loạt sững lại: đôi mắt ấy\dots\ không phải xanh thẳm như Shino mà chúng tôi quen biết, mà là một nâu dịu, hoàn toàn khác biệt.

Người anh trai nuốt khan, giọng gần như thì thầm với chính mình: ``\dots Sora-san.''

Saikyou nheo mắt, đôi môi mím chặt, rồi khẽ gật, xác nhận điều cả hai vừa nhận ra. ``Đôi mắt nâu kia\dots\ không thể lẫn vào đâu được.''

Tôi vẫn chưa hiểu hết. Tôi nghe rõ cái tên họ nói ra nhưng không thể tin nổi. ``Sora\dots?'' Tôi lập lại, quay trở về nhìn gương mặt người con gái trước mặt.

Đúng lúc ấy, người đó mỉm cười gượng gạo, ngập ngừng lên tiếng: ``Phải\dots\ Tớ là Sora.''

Trái tim tôi chao đảo. Trong đầu vẫn còn vọng lên cái tên quen thuộc mà miệng tôi vừa bật ra trong vô thức. ``S-Shino-san\dots\ Chuyện gì\dots\ xảy ra với cậu vậy\dots?''

Tôi vẫn không thể tin được người trước mặt đó là Sora. Rõ ràng dáng vóc đó, cái cài tóc đó, và cả giọng nói nhẹ nhàng kia\dots\ tôi không hề nhầm. Rõ ràng đó là Shino-san, không ai nhầm lẫn được.

``Shino-san!'' tôi gọi lớn, giọng run rẩy. ``Cậu sao vậy? Tại sao lại nói những điều kỳ lạ thế?''

Nhưng cô gái trước mặt chỉ lắc đầu, ánh mắt vẫn giữ nguyên vẻ bình thản đến kỳ lạ. ``Tớ nghĩ cậu nhận nhầm ai đó rồi. Tớ\dots\ thực sự là Sora mà.''

``Không thể nào\dots'' tôi lẩm bẩm, cảm giác hoang mang càng lúc càng dâng cao. ``Chắc chắn là do cú sốc vừa rồi\dots\ Phải rồi, cậu bị mất trí nhớ\dots\ Hay là\dots''

``Reiko-san,'' giọng Saikyou-san đột ngột cắt ngang dòng suy nghĩ của tôi. ``Cô ấy nói đúng. Người trước mặt cậu không phải Shino-san đâu.''

Tôi quay phắt lại, nhìn cô ấy với ánh mắt không tin. ``Sao cậu lại nói vậy? Rõ ràng là\dots''

``Chúng tôi đã gặp cô ấy trước đây,'' Kyuen-san chậm rãi lên tiếng, giọng trầm và chắc chắn, dù chính cậu cũng không thể tin nổi. ``Trong giấc mơ, và cả khi chúng tôi tìm đường thoát khỏi trường học sau trận động đất ở thời đại của chúng tôi.''

Cô em gái gật đầu, dường như không cần nói thêm câu nào nữa.

Tôi đứng chết lặng, nhìn qua nhìn lại giữa hai anh em họ và người con gái tự xưng là Sora. Trong khi tôi vẫn bám víu vào ý nghĩ đây là Shino-san bị ảnh hưởng tâm lý, thì cả Saikyou-san và Kyuen-san đều tỏ ra bình tĩnh đến kỳ lạ, như thể họ đã quen với tình huống này.

``Nhưng\dots\ làm sao có thể\dots'' tôi lắp bắp, cố gắng tìm lời giải thích hợp lý. ``Làm sao mà\dots''

``Reiko-san,'' giọng Sora nhẹ nhàng vang lên. ``Tớ hiểu cảm giác của cậu. Nhưng tin tớ đi, tớ không phải là người cậu đang tìm kiếm.''

Tôi nhìn kỹ hơn vào gương mặt cô ấy. Đúng là có điều gì đó rất khác - không đơn thuần chỉ ở hình dáng bên ngoài, mà là thứ gì đó sâu hơn, cơ bản hơn. Cách cô ấy nhìn, cách cô ấy nói chuyện, tất cả đều mang một sắc thái hoàn toàn khác với Shino-san mà tôi quen biết.

Tôi vẫn cố tìm một lời giải thích hợp lý cho sự khác biệt này, nhưng cổ họng ứ lại, chỉ còn lại những câu hỏi chồng chất.

``\dots Nếu cậu thật sự không phải Shino-san, vậy\dots\ cậu biết gì về những người khác không?'' tôi buột miệng hỏi, như để tìm bằng chứng phủ nhận chính lời cô ấy vừa nói.

Sora-san - theo cô ấy tự xưng - khẽ nghiêng đầu, ánh mắt nâu trầm như đang gợi nhớ lại. Cô mỉm cười, dịu dàng nhưng cũng có đôi chút nặng nề. ``Nanase-san\dots\ cậu ấy luôn cứng cỏi, thẳng thắn, chẳng bao giờ chịu tha thứ cho những kẻ giả dối. Đôi khi hơi gắt, nhưng thực lòng thì chỉ muốn bảo vệ mọi người.''

Tôi khẽ giật mình. Đúng là Nanase-san vẫn hay nổi cáu vì những chuyện nhỏ nhặt, đặc biệt là khi ai đó nói dối.

Cô tiếp tục, giọng đều đều, nhưng ánh mắt như lạc về một nơi xa xăm: ``Shiina-san thì khác. Cậu ấy ồn ào, cá biệt, thích gây gổ, nhưng thật ra lại để ý đến từng chi tiết xung quanh. Cậu ấy nhìn mọi thứ rất kỹ, kể cả những điều người khác bỏ qua.''

Tôi thoáng ngẩn người. Những lần Shiina-san mắng tôi vì bất cẩn thoáng chốc ùa về. Saikyou-san cũng khẽ gật đầu, còn Kyuen-san thì mím môi đồng tình.

``Honoka-san thì khỏi phải nói,'' Sora-san cười khẽ, ``lúc nào cũng là người pha trò, luôn nghĩ ra những trò khiến bầu không khí nặng nề thành nhẹ nhõm chỉ bằng vài câu bông đùa. Cậu ấy cũng là người bày ra những ý tưởng khó đỡ.''

Tôi bỗng nhớ lại cảnh cô bạn tóc vàng nhảy nhót trong lớp, bắt chước giọng giáo viên chỉ để chúng tôi phá lên cười. Chính cậu ta là người đề ra thử thách gan dạn chết chóc ấy. Một thoáng buốt nhói len lỏi trong tim.

``Còn Yae-san\dots\ cô ấy siêng thể thao lắm, luôn lôi kéo mọi người cùng tập chạy, và hay giảng giải về việc vẻ đẹp phụ nữ truyền thống, và việc họ cần tự lập, không nên phụ thuộc.''

Tôi nghe mà không thốt nên lời. Những buổi sáng chạy thể dục cùng Yae-san hiện về rõ ràng, cả cái cách cô ấy luôn nói triết lý về phụ nữ.

``Sakura-san thì\dots'' Sora ngập ngừng một thoáng, rồi khẽ thở dài. ``Cô ấy thường hay trốn học, lúc nào cũng viện đủ lý do. Nhưng khi nói chuyện, cô ấy lại thích tranh luận tới cùng, chẳng bao giờ chịu thua. Cậu ấy cũng là một người thông minh.''

Tôi nín lặng. Đúng từng chữ. Đó là Sakura-san, không lẫn đi đâu được.

Tôi quay sang nhìn hai anh em họ. Cả hai đều im lặng, nhưng gương mặt họ cho thấy rõ rằng những gì Sora vừa kể không hề sai.

Tôi mím môi, lồng ngực như bị ai đó bóp chặt. Cô ấy\dots\ biết tất cả. Từng người một. Nhưng tại sao\dots\

``C-Còn mình thì sao?'' tôi hỏi, giọng run rẩy, tim đập thình thịch. ``Cậu biết gì về mình?''

Sora khựng lại. Đôi mắt nâu ấy nhìn tôi, thoáng bối rối. Rồi cô khẽ lắc đầu. ``Xin lỗi\dots\ nhưng tớ chưa từng gặp cậu bao giờ.''

Lời đáp như một nhát dao xoáy vào lồng ngực tôi. Tôi lùi lại nửa bước, không tin nổi vào tai mình. Rõ ràng cô ấy nhớ hết mọi người, từng chi tiết nhỏ\dots\ trừ tôi.

``Không\dots\ không thể nào\dots'' tôi thì thào.

Sora vẫn giữ ánh mắt dịu dàng, nhưng ẩn chứa nỗi áy náy khó tả. Kyuen và Saikyou nhìn nhau, rồi nhìn tôi, nhưng chẳng ai mở miệng.

``Vậy\dots\ nếu cậu nói cậu là Sora\dots\ thì Shino-san thật sự đang ở đâu?'' tôi hỏi, giọng run rẩy.

Một khoảng lặng nặng nề bao trùm. Sora nhìn tôi với ánh mắt trắc ẩn, trong khi hai anh em họ đưa những cái nhìn lo lắng lẫn nhau, không biết phải giải thích như thế nào để cho tôi hiểu.

Thực sự\dots\ chuyện gì đang xảy ra vậy? Tại sao Shino-san không có ở đây? Tại sao chỉ có Sora-san? Có phải đây là trò đùa của thế lực siêu nhiên nào đó? Hay đây là một cơn ác mộng mà tôi không thể tỉnh dậy? Tại sao mọi thứ lại trở nên kỳ lạ và đáng sợ như thế này?

Hàng tấn câu hỏi cứ vậy, nhưng không có một câu trả lời thỏa đáng nào.

\pov{-}

Không gian quanh họ dường như nén lại, u ám hơn bao giờ hết. Giữa đống đất đá, gạch vụn và những thanh gỗ gãy nát vắt chéo nhau, khung cảnh chẳng khác gì tàn tích của một bãi chiến trường vừa bị càn quét. Không còn dấu vết nào cho thấy nơi đây từng là một ngôi trường, nơi mà các học sinh cắp sách tới trường, nơi tri thức được trao, mà chỉ còn là đống ngổn ngang sau một cơn giận dữ của thiên nhiên.

\img{e1-8b}

Sora đứng lặng người, đôi mắt nâu đảo khắp cảnh hoang tàn. Giọng cô khẽ run: ``Chuyện\dots\ chuyện gì đã xảy ra thế này?''

Không ai trả lời. Chỉ có tiếng gió thổi qua những khe nứt, và tiếng bụi rơi lả tả từ những mảng tường đổ. Sora siết chặt nắm tay, hỏi lại: ``Làm ơn nói cho tớ biết\dots\ đây là chuyện gì?''

Im lặng vẫn bao trùm. Kyuen và Saikyou cúi đầu, không dám nhìn thẳng vào mắt cô. Sora cảm thấy một nỗi sợ hãi dâng lên trong lồng ngực, giọng cô run rẩy hơn: ``Xin các cậu\dots\ nói gì đi\dots''

Cuối cùng, Reiko chua chát trả lời, giọng lạc hẳn đi: đó là một trận sạt lở đất, đã chôn vùi toàn bộ ngôi trường chỉ trong chớp mắt. Sora nghe, và gương mặt lập tức biến sắc. Trong ký ức mơ hồ và vụn vỡ vốn thuộc về Shino, cô biết rằng điều này từng xảy ra. Song sự thật lại phũ phàng giáng xuống, khiến đôi vai cô khẽ run.

Ánh mắt Sora trượt sang Kyuen và Saikyou. Cô em gái đang gồng mình cõng ông anh trai thương tích đầy mình, đầu bê bết máu, sắc mặt trắng bệch. Dù Saikyou cố gắng giữ bình tĩnh, từng bước chân vẫn chao đảo, cho thấy cô đang phải kiệt sức thế nào.

Sora vội bước đến gần, đôi bàn tay run rẩy chạm nhẹ vào vết thương trên đầu Kyuen.

``Cậu có sao không? Vết thương sâu lắm! Chúng ta phải băng bó ngay!'' cô hỏi dồn dập, giọng lạc đi vì lo lắng, chẳng còn giữ được vẻ xa cách hay ngập ngừng ban đầu.

``Không sao đâu\dots\ tôi vẫn chịu đựng được mà,'' Kyuen miễn cưỡng đáp, nhưng sự thật quá rõ ràng: cơ thể cậu chẳng thể gắng thêm nữa.

``Và chính vì thế nên em mới phải cõng anh đấy, Nii-chan,'' cô em gái của cậu nhíu mày.

Khoảnh khắc yên lặng trĩu nặng tràn đến. Sora ngập ngừng cất tiếng, giọng mong manh như chỉ còn là hơi thở: ``Vậy\dots\ mọi người thì sao? Họ\dots\ họ vẫn ổn cả chứ?''

Không ai trả lời ngay. Sự im lặng tiếp tục bao trùm không gian, nặng nề như những tảng đá vừa sập xuống. Sora lên tiếng, giọng run rẩy vì sự im bặt đáng sợ kia: ``Các cậu\dots\ tìm thấy họ chưa?''

``Yae-san\dots\ chết rồi,'' Kyuen nói, giọng trầm xuống.

``Shiina-san cũng vậy,'' Saikyou tiếp lời, nước mắt lăn dài trên má.

``Honoka-san và Nanase-san\dots\. cũng không còn nữa,'' Reiko nghẹn ngào, đôi tay run rẩy.

Sora chết lặng. Đôi mắt nâu ấy như vừa bị dập tắt, ánh sáng bên trong vụt tàn. Hơi thở trở nên gấp gáp, ngắt quãng.

``V-Vậy còn Sakura-san?'' cô hỏi, giọng run rẩy.

``Có lẽ cũng không còn hy vọng nữa. Bọn tôi không thể tìm thấy cô ấy,'' Kyuen đáp. Kèm theo một cái lắc đầu nhẹ của cô em gái.

Sora chết lặng. Đôi mắt nâu ấy như vừa bị dập tắt, ánh sáng bên trong vụt tàn. Hơi thở trở nên gấp gáp, ngắt quãng. Cô khẽ lẩm bẩm, như đang tự nói với chính mình.

``Không phải\dots\ lại nữa chứ\dots\ Mình\dots\ lại không thể cứu họ rồi\dots''

Cô lảo đảo đứng dậy. Đôi giày lười phủ bụi nặng nề lê bước qua những tảng gạch vỡ. Gió lạnh thổi qua những khe nứt trên tường, tạo thành những tiếng rít ghê rợn như tiếng khóc than. Bụi từ trần nhà rơi lả tả xuống vai cô, như một cơn mưa tro tàn không dứt. Không gian chìm trong một sự im lìm đến đáng sợ, chỉ còn tiếng vữa vụn rơi lộp độp và tiếng gió hú qua những ô cửa trống hoác. Không một lời giải thích, không một ánh nhìn ngoái lại, Sora chỉ lặng lẽ rời đi khỏi đống tàn tích, bóng dáng gầy guộc hòa vào nền xám lạnh của hoang tàn, như một linh hồn lạc lõng giữa những đổ nát của thời gian.

``\textbf{Khoan đã! Sora-san!}'' Reiko kêu lên, cùng với tiếng gọi của Saikyou. Cả ba vội vã đuổi theo, nhưng đáp lại họ chỉ là sự im lặng tuyệt đối. Sora bước đi, kiên định đến lạ, như thể bị dẫn dắt bởi một sức mạnh vô hình nào đó.

Không còn lựa chọn nào khác, cả ba đành phải nối gót theo sau, trái tim nặng trĩu bởi những mất mát chưa kịp nguôi, và một nỗi bất an mơ hồ đang chờ đợi phía trước.

\ct

Sora bước đi loạng choạng, mái tóc dài khẽ lay động trong làn gió âm u phả qua từ phía núi. Đằng sau, Kyuen được Saikyou cõng trên lưng, hơi thở nặng nề, máu khô loang thẫm bên thái dương. Reiko, mặt mày lấm lem bụi đất đi sau cùng. Cả bốn rời xa đống hoang tàn, bỏ lại phía sau nơi từng gọi là ``ngôi trường'', giờ chỉ còn là đống gạch vụn không hơn không kém.

Con đường dẫn lên núi thoai thoải, lát những bậc đá đã sứt mẻ theo năm tháng. Trên cao, sương mù hơi phủ, khiến cảnh vật càng thêm huyền hoặc dưới ánh trăng tròn. Vầng trăng sáng vằng vặc treo lơ lửng giữa bầu trời đêm trong vắt, tỏa ánh bạc xuống mặt đất.

Bốn người bước đi, bóng họ đổ dài trên những bậc thang đá, lung linh dưới ánh trăng như những bóng ma vật vờ giữa màn đêm tĩnh mịch. Ánh trăng phủ lên vai họ một lớp áo bạc mỏng manh, dáng người mờ ảo.

\img{e1-8c}

Một cánh cổng torii đỏ sẫm hiện ra trước mặt. Sora đi xuyên qua, không ngẩng đầu, còn ba người kia dừng lại chốc lát, như chợt thấy điều gì kỳ lạ. Người con trai khẽ nhăn mặt, cất giọng khàn: ``Chúng ta\dots\ đã từng đi qua chỗ như thế này chưa?''

Cô em gái siết chặt tay cõng cậu, thì thầm đáp, ``Không, em không nhớ. Nhưng nơi này\dots\ quá yên tĩnh\dots''

Cô gái mắt hai màu chỉ lặng lẽ đưa mắt nhìn cánh cổng đỏ, lòng dâng lên nỗi bất an mơ hồ. Chỏm tóc ngố của cô gập xuống, linh cảm chẳng lành.

Họ tiếp tục tiến lên. Xen kẽ giữa các cổng torii đỏ ấy là những bụi hoa bỉ ngạn mọc san sát hai bên đường. Những cánh hoa đỏ rực, rũ xuống như những đốm lửa đang cháy âm ỉ, tương phản sâu sắc với sắc xám của trời đất. Đây không phải lần đầu cả ba nhìn thấy loài hoa này, chúng từng hiện hữu nhiều lần trong quá khứ, và cả trong cơn ác mộng họ đã từng trải, nhưng như mọi lần, nó là một điềm báo chẳng lành.

\img{e1-8d}

Nhưng vì sao chúng lại mọc trên đoạn đường này, nơi vốn không nên tồn tại thứ gì ngoài đá và rêu, thì chẳng ai có lời giải thích. Một thế lực siêu nhân nào đó? Hay đây là mùa của hoa? Không ai biết được.

Saikyou chép miệng khe khẽ, ánh mắt né tránh: ``Không đúng\dots\ chỗ này\dots\ đáng lẽ không có hoa nào cả.''

Reiko rùng mình, nắm chặt vạt áo, lẩm bẩm: ``Tại sao chúng lại ở đây được nhỉ\dots?''

Cứ thế, họ bước qua những cổng torii đỏ nối tiếp nhau, như đang lạc vào một hành lang vô tận giữa hai thế giới. Cuối cùng, khi màn sương mỏng tản dần, ngôi đền hiện ra.

Ngôi đền cổ bao quanh bởi khu rừng rậm rạp, thân cây cao lớn vươn lên như những bóng đen sừng sững. Bậc đá dẫn vào đã rạn nứt, rêu xanh bám đầy. Hai chiếc đèn đá lặng lẽ đứng chầu bên lối đi, phủ bụi mờ. Mái đền cong vút, chạm khắc những hoa văn cổ xưa mờ nhạt, từng đường nét như đã nhuốm bụi thời gian. Từng mảnh gỗ cũ kêu lên cọt kẹt dưới bước chân, nhưng ngôi đền vẫn giữ vẻ uy nghiêm lạnh lẽo. Nói cách khác, giống hệt nơi mà Shino từng đến cầu nguyện.

\begin{figure}[H]
  \centering
  \includegraphics[width=12.5cm]{../../../../Image/Book?/e1-6a.png}
  \caption{Về cơ bản, ngôi đền mà Shino giống hệt như bức ảnh này, nhưng diễn ra vào ban đêm và không có mưa.}
\end{figure}

Saikyou dừng lại trước cổng đền, hơi thở gấp gáp vì đã không còn sức lực. Kyuen nghiêng đầu nhìn vào, đôi mắt vẫn còn mơ hồ vì cơn đau. Reiko đứng sát bên, lòng nặng trĩu. Không ai dám bước vào.

Chỉ có Sora là lặng lẽ tiến về phía trước, dáng người nhỏ bé nhưng kiên định. Cô bước qua ngưỡng cửa gỗ, đi thẳng vào bên trong, không ngoái lại.

Ba người còn lại chỉ biết đứng ngoài, nhìn bóng dáng ấy tan vào bóng tối của ngôi đền.

``Sora-san\dots'' Saikyou-san khẽ lên tiếng, giọng lo lắng. ``Cậu cấy liệu có ổn không?''

Reiko liếc nhìn cô gái tóc nâu đang lặng lẽ bước đi phía trước. Từ khi rời khỏi nơi ấy, Sora không nói một lời nào, chỉ im lặng bước đi như một cái bóng vô hồn. Dáng đi của cô có gì đó rất khác - không còn vẻ mạnh mẽ thường ngày, mà như thể mọi sức sống đã bị rút cạn.

``Cô ấy\dots\ chắc đang rất đau lòng,'' Kyuen thì thầm, giọng trầm xuống. ``Dù sao thì\dots\ Sakura-san cũng là người bạn thân nhất của cô ấy.''

Reiko gật đầu, cảm thấy một nỗi buồn sâu thẳm dâng lên trong lòng. Họ đều hiểu cảm giác mất đi người thân thiết nhất với mình là như thế nào. Nhưng với Sora, có lẽ nỗi đau ấy còn sâu sắc hơn gấp bội.

``Chúng ta nên\dots'' cô ngập ngừng, ``nên làm gì đó không?''

Saikyou lắc đầu nhẹ. ``Đôi khi\dots\ im lặng là cách tốt nhất để chia sẻ nỗi đau.''

Cả ba người tiếp tục bước đi trong yên lặng, mắt không rời khỏi dáng hình cô độc phía trước. Dù không ai nói ra, nhưng tất cả đều hiểu rằng: đằng sau vẻ ngoài lạnh lùng kia là một trái tim đang tan vỡ từng mảnh.

Theo đúng nghi thức, Sora nắm lấy sợi dây thừng lớn treo trước gian chính, lắc mạnh. Tiếng chuông đồng ngân dài, lan ra khắp không gian âm u như thể xuyên qua cả núi rừng. Cô bỏ vài đồng xu vào hòm công đức, rồi cúi đầu hai lần, vỗ tay hai cái. Đôi mắt nâu ngấn lệ khép lại, đôi môi run rẩy bật thành lời cầu nguyện.

``Đây\dots\ có lẽ là một cơn ác mộng. Điều này\dots\ không thể xảy ra.''

Khoảng lặng phủ xuống. Bên ngoài, Reiko nắm chặt vạt áo, khẽ lắc đầu khi nhớ lại cảnh tượng kinh hoàng ấy. Kyuen nhíu mày, đôi mắt mờ đục vì vết thương nhưng vẫn dõi theo chăm chú. Saikyou chỉ cắn môi, không nói gì.

Bên trong, Sora khẽ run, nước mắt bắt đầu rơi trên má. ``Thần linh ơi\dots\ Ngài đã cho con kết bạn với nhiều người như vậy\dots\ để rồi\dots'' giọng cô nghẹn lại, bàn tay siết chặt đến tái nhợt.

\img{e1-8e}

Reiko hít sâu, thì thầm: ``Sora-chan\dots'' nhưng rồi lại im lặng, không dám gọi to. Đôi mắt cô cũng lặng lẽ rơi lệ trước một cô gái đã sụp đổ, dù cô mới chỉ gặp cách đó không lâu.

``\textbf{Con xin Thần linh đấy\dots\ Hãy cứu họ đi!}'' tiếng kêu bật ra, lạc đi vì nức nở, hòa lẫn tiếng chuông ngân đang dần tắt trong gió. Nước mắt lăn dài, nhỏ xuống sàn gỗ lạnh, từng giọt nặng trĩu.

Kyuen khẽ rên một tiếng vì đau, nhưng ánh mắt anh vẫn như dõi theo ngọn lửa yếu ớt còn sót trong lời cầu nguyện của cô gái ấy.

``\textbf{Con cầu xin Thần linh hãy thỉnh nguyện lời cầu của con!}''

Vừa dứt lời, tiếng chuông vang lên lần nữa, nhưng lần này mạnh mẽ đến lạ. Toàn bộ không gian rung chuyển, ánh sáng trắng bất chợt bừng lên từ trong gian chính, soi rọi từng góc đền, rồi lan rộng ra ngoài.

\img{e1-8f}

``\textbf{Chói quá!}'' Reiko kêu lên, hai tay ôm mặt.

``\textbf{Mắt em!}'' Saikyou hét, giọng hoảng loạn. ``\textbf{Em không thấy gì cả!}''

``\textbf{Đừng buông tay nhau ra!}'' Kyuen gào lên, cố với lấy em gái.

Reiko giật mình che mắt. Saikyou khụy gối xuống, nhưng vẫn gắng cõng anh trai. Kyuen cố nhướng mắt, ngạc nhiên lẫn bất an.

Họ chưa kịp làm gì thêm thì ánh sáng ấy lan tràn, trắng lóa như xóa nhòa mọi hình thù, mọi màu sắc, nuốt chửng mọi thứ.

``Chuyện gì đang\dots\ xảy ra vậy?'' giọng Reiko yếu dần.

``Anh\dots\ em không cảm thấy gì cả\dots'' Saikyou thì thầm, như đang chìm vào giấc ngủ.

``Đừng\dots\ đừng để em một mình\dots'' Kyuen lẩm bẩm trước khi giọng nói tan biến.

Ngôi đền, cánh rừng, con đường với hoa bỉ ngạn, cả bốn người, tất cả đều bị bao trùm trong ánh sáng ngợp trời. Cảnh vật rung động rồi biến mất, dần dần tan chảy trong một màu trắng vô tận.

\dots

\dots

\ct

\pov{Haragen Kyuen}

\dots

\dots

``Ư\dots\ chuyện quái gì vừa xảy ra vậy\dots?''

Ánh sáng trắng lóa dần dịu xuống. Phía trước mặt tôi giờ là một màu đen kịt, chẳng nhìn thấy một cái gì cả. Chỉ có bóng tối bao xung quanh.

Một cơn nhức đầu nhẹ thoáng qua tôi, nhưng tôi không cảm thấy khó chịu. Nó ập đến từ từ, rồi dần dịu đi, giống như tôi vừa mới ngủ dậy.

Nhắc tới ngủ dậy\dots\ tôi nhận ra mình không còn đứng nữa, mà đang nằm trên một cái gì đó\dots\ hoặc đúng hơn, là trên ai đó. Cơ thể tôi hơi nặng, lồm cồm tìm điểm tựa để chống dậy.

``Tại sao\dots\ mình lại nằm thế này\dots?''

Nhưng ngay khi hai bàn tay tôi chạm xuống, một cảm giác mềm và căng làm tôi giật mình.

``Ahn\~{}'' ``Eek!'' hai tiếng kêu khẽ bật ra cùng lúc từ hai phía.

Tôi quay sang bên trái, rồi bên phải. Khung cảnh phía trước tôi làm tôi khựng lại: tay trái đang ép xuống ngực của Saikyou, còn tay phải thì đặt hẳn lên bụng Reiko-san. Cả hai vừa tỉnh lại, đôi má ửng lên vì bất ngờ.

``Nii-chan\dots\ đây không phải chỗ để anh làm vậy đâu\dots'' em tôi tròn mắt, giọng có một chút nổi cáu xen lẫn ngượng ngùng. Đến tận đây rồi em ấy vẫn còn đùa được, tài thật\dots

Reiko-san thì che vội lấy phần bụng, nhưng ánh mắt lại chuyển ngay sang lo lắng. ``Cậu\dots\ ổn chứ? Vết thương trên đầu cậu\dots''

Tôi đưa tay sờ lên thái dương mình. Không còn cảm giác đau rát hay máu chảy nữa. Thay vào đó, da thịt lành lặn, như chưa từng bị thương. Thân thể cũng nhẹ nhõm, khỏe khoắn như trước khi mọi chuyện xảy ra. Một phép lạ gì đó từ bên trong ánh sáng trắng đó ư?

Saikyou lập tức bật dậy, nhìn chằm chằm tôi, rồi khẽ thở phào: ``Quả thật\dots\ vết thương của anh đã biến mất rồi.''

``Hai người có sao không?'' tôi hỏi, giọng lo lắng.

``Bọn em vẫn ổn\dots'' cả hai người đáp lại.

``Thật sự không sao chứ? Có chắc không bị thương gì không?''

``Hoàn toàn ổn mà! Nhìn này!'' Saikyou nói rồi dang rộng hai tay, xoay một vòng tròn để chứng minh. Ngoài mái tóc hơi rối, em ấy trông vẫn khỏe mạnh như thường.

``Mình thề đấy, không có gì đáng ngại cả!'' Reiko-san khẳng định, giọng chắc nịch. ``Cậu lo xa quá rồi.''

Nhưng dường như vẫn chưa đủ thuyết phục, em gái tôi bỗng nảy ra một ý: ``Nếu anh vẫn không tin thì\dots\ em có thể cho anh kiểm tra kỹ hơn\dots''

``S-Saikyou-san!? Khoan đã-- Á!?''

Không đợi cô ấy kịp phản ứng, em tôi đã nhanh tay túm lấy gấu váy của cô bạn, định kéo lên. Reiko-san hốt hoảng lùi lại, mặt đỏ bừng, trong khi Saikyou vẫn cười tinh nghịch, cố với tay về phía trước. ``Để em chứng minh cho anh thấy là bọn em không sao thật mà\~{}''

Tôi lập tức giơ hai tay lên, mặt nóng ran. ``T-Thôi! Không cần! \textbf{Anh tin rồi!! Tin rồi mà!!!}''

Trời ạ, em tôi lúc nào cũng vậy\dots\ Hình như tôi đã từng gặp chuyện này ở đâu đó rồi thì phải\dots

``Cơ mà nhé, hai người\dots'' Reiko-san bất chợt lên tiếng, hoàn hồn sau vụ vừa rồi.

``Chúng ta\dots\ đang ở đâu đây?'' tôi hỏi, nhìn quanh bối rối.

``Rõ ràng ta vừa ở ngôi đền mà\dots'' Saikyou lẩm bẩm.

Chúng tôi mới bắt đầu nhận ra xung quanh. Tất cả đang ở trong một hành lang quen thuộc: những ô cửa kính, hàng ghế gỗ, mùi sơn cũ kỹ. Reiko-san chậm rãi thì thầm: ``Đây\dots\ là hành lang của trường Aogami. Nhưng rõ ràng\dots\ nó đã sập hoàn toàn rồi mà\dots''

Tôi nhìn quanh, gật đầu chậm rãi. ``Ừ\dots\ nhưng có gì đó không đúng.''

Saikyou tiếp lời, giọng nhỏ đi, khẽ gật: ``Nơi này\dots\ giống nhưng không phải. Giống như là một bản sao nào đó.''

``Kiểu như là\dots\ ở trong giấc mơ sao?'' cô nàng mắt hai màu hỏi.

``Dạng như vậy.''

Ngay trước mặt, một chiếc đèn pin nằm chỏng chơ. Reiko-san cúi xuống, nhặt lên, bấm thử. Ánh sáng trắng vàng hắt ra, rọi xuống sàn gỗ bóng ẩm, kéo theo những vệt sáng dài ngoằn ngoèo, nhưng không đủ để rọi tới bên kia.

Chúng tôi bắt đầu đi, từng bước chậm rãi. Người cầm chiếc đèn pin đi trước, trong khi hai anh em tôi ngay sát theo sau, trong lòng không giấu được sự bất an, dù chúng tôi đã từng rơi vào tình cảnh này không ít lần.

Tiếng bước chân vang vọng khắp hành lang, giữa một không gian vừa quen vừa lạ, vừa như của quá khứ, vừa như đã bị biến dạng bởi một sức mạnh gì đó vô hình.

Không khí trong hành lang dày đặc đến mức ngực tôi thấy khó thở. Mọi âm thanh biến mất, chỉ còn lại tiếng ù ù văng vẳng như vọng từ lòng đất. Đèn pin trong tay Reiko-san cắt ngang màn tối, nhưng ánh sáng ấy chỉ đủ soi một khoảng nhỏ, làm cái bóng của chúng tôi in dài ngoằn ngoèo lên tường.

``Reiko-san\dots\ nơi này có vẻ im ắng thật,'' tôi thì thầm, cảm nhận được sức nặng của Saikyou đang dựa hẳn vào người mình.

``Nơi này khiến cho mình cảm thấy có gì đó thân thuộc\dots'' Reiko-san đáp, giọng lo lắng.

``Tốt nhất là cứ cẩn thận thì hơn\dots'' em tôi lên tiếng, nhỏ giọng.

Chúng tôi đi dọc theo hành lang im ắng, và rồi người đi trước mặt bỗng khựng lại.

``T-Tiếng gì vậy\dots?'' Saikyou thì thầm, giọng run rẩy.

Cô lia ánh sáng về phía một lớp học vẫn còng sáng đèn bên trong. Tim tôi đập nhanh hơn. Cả ba nghiêng người, nhìn qua khe cửa.

Đó là một lớp học. Trông nó mang dáng dấp của lớp 1-B chúng tôi. Nơi giờ đã bị vùi trong đống đổ nát, vậy mà ở đây nó lại hiện ra, sáng sủa, nguyên vẹn. Bên trong, những bóng đen mang dáng dấp con người chuyển động qua lại. Tôi nhận ra nét quen thuộc: hình hài của những người bạn mà Reiko đã mất.

``Không thể nào\dots'' Reiko-san thở gấp, bàn tay cầm đèn pin run lên. ``Họ\dots\ họ đã\dots''

``Đừng nhìn nữa,'' tôi kéo tay cô ấy lại. ``Đó không phải họ đâu.''

``Ta phải chấp nhận sự thật thôi, Reiko-san\dots'' em tôi vỗ vai cô ấy.

Những bóng đen đó\dots\ Không chỉ là học sinh từ năm 1946, mà còn là những người đã từng học ở đây trong thời hiện đại nữa. Tôi và Saikyou đã thấy họ trong giấc mơ, những linh hồn vất vưởng không thể siêu thoát.

Và giờ, chúng lại ở đây, dù mang hình hài khác, bản chất vẫn như vậy.

Chúng đang nói cười, trò chuyện rôm rả. Âm thanh như tiếng cười học sinh, nhưng méo mó, vang vọng như bị kéo dài, ma mị. Không khí lạnh lẽo xuyên qua sống lưng tôi.

*XẸT XẸT*

``C-Cái gì vậy!?'' người cầm chiếc đèn pin giật thót, lùi lại về sau.

Đột nhiên, những bóng đen kia kéo giãn, hình thù méo mó như bị lỗi, rồi vụt biến mất. Căn phòng sáng đèn lại trống rỗng, tĩnh mịch đến rợn người.

``Chúng biến mất rồi\dots'' tôi thì thầm, vẫn chưa hết bàng hoàng.

``Kyuen-san, những thứ đó là gì vậy?'' Reiko-san hỏi, giọng run run.

Tôi và Saikyou dẫu cũng giật mình nhưng vẫn vội trấn tĩnh cô. Ánh mắt cô run rẩy, hoang mang, như thể đã từng nhìn thấy cảnh tượng ấy trong một giấc mơ.

``Hai anh em tôi cũng từng thấy cảnh này,'' Saikyou nói khẽ, ``trong những cơn ác mộng\dots\ lúc mà chúng tôi tìm đường rời khỏi hành lang trường.''

``Hình như\dots\ mình đã từng thấy chúng trong những giấc mơ như vậy,'' Reiko-san đáp, tay lẩy bẩy cầm chiếc đèn pin.

Tôi nhớ ra, Reiko từng kể qua về giấc mơ đó\dots\ và những sự trùng hợp của chúng làm tôi lạnh gáy.

``Chúng ta phải cẩn thận,'' tôi nói, ``những cái bóng đó không bình thường.''

Cô gái hai màu mắt không còn cách nào khác ngoài việc gật đầu. Chúng tôi tiếp tục tiến về phía trước. Hành lang dẫn tới một ngã ba, và trước mặt là một cánh cửa khép kín. Saikyou bước lên, định với tay mở. Nhưng đúng lúc ấy\dots\

*CẠCH!*

\dots một âm thanh vang lên từ hành lang bên phải.

``L-Lại gì nữa vậy!?'' Reiko-san hét lên, giọng run rẩy. Cô vội vàng lia đèn pin về phía phát ra tiếng động, tay siết chặt thiết bị đến mức các khớp ngón tay trắng bệch.

Ánh sáng đèn pin chạm vào một bóng đèn huỳnh quang treo lơ lửng trên trần, đang chập chờn một cách kỳ quái. Dưới ánh sáng nhấp nháy ấy, một chiếc ghế gỗ đơn độc hiện ra. Trên ghế là một bóng người ngồi im lìm, bất động như một pho tượng.

``C-Có người!?'' cô ấy hoảng lên, vô thức nói lớn, nhưng\dots\ bóng người ấy không có phản ứng gì cả.

``Cẩn thận,'' tôi thì thầm, giọng căng thẳng. ``Đừng lại gần quá.''

``Em có cảm giác không ổn về chuyện này\dots'' em tôi nắm chặt tay, giọng lo lắng.

Chúng tôi từ từ tiến lại, mỗi bước chân đều dè dặt, thận trọng. Tim tôi đập thình thịch trong lồng ngực. Mỗi khi ánh đèn huỳnh quang lóe lên rồi tắt phụt, bóng người trên ghế như méo mó, biến dạng theo từng nhịp sáng tối.

``C-Có ai ở đó không?'' Reiko-san cất tiếng hỏi, nhưng chỉ có sự im lặng đáp lại.

Khi chúng tôi chỉ còn cách vài bước, bóng đèn đột ngột tắt phụt, nhấn chìm tất cả trong bóng tối dày đặc.

``Kyaa!'' Saikyou-san thét lên, bấu chặt vào tay áo anh trai.

``Suỵt! Bình tĩnh nào\dots'' tôi đặt tay lên vai em gái. ``Nghe kìa\dots''

Trong khoảnh khắc tối đen ấy, một âm thanh kỳ lạ vang lên - như tiếng gì đó di chuyển, trượt trên sàn gỗ cũ kỹ.

``C-Cái gì vậy!?'' Reiko-san run giọng, vội chiếu đèn pin về phía phát ra âm thanh.

Khi ánh sáng chập chờn của bóng đèn huỳnh quang trở lại, cả chiếc ghế và bóng người đã biến mất không dấu vết. Thay vào đó, một vũng chất lỏng đỏ thẫm đang từ từ lan rộng dưới chân, mùi tanh nồng của máu xộc thẳng vào mũi chúng tôi.

``Máu\dots\ Máu kìa\dots' Saikyou thì thầm, mặt tái nhợt. ``Đó\dots\ là máu phải không?''

``Không thể nào\dots'' tôi lẩm bẩm, giọng nghẹn lại. ``Làm sao có thể\dots''

Reiko-san lùi lại một bước: ``Chuyện quái gì đang xảy ra vậy\dots''

Không khí xung quanh chúng tôi nặng nề và ngột ngạt. Mùi máu tanh tưởi cùng với lay lắt hạt bụi tạo nên một hỗn hợp khó thở. Thỉnh thoảng, một cơn gió lạnh thổi qua khe nứt, mang theo tiếng rên rỉ của những thanh gỗ đang oằn mình chống đỡ. Ánh sáng yếu ớt chiếc đèn pin tạo nên những vệt sáng mờ ảo, chiếu rọi lên khung cảnh trước mắt. Cảm giác như thể chúng tôi đang bị nhốt trong nhà lao vậy.

``Chúng ta phải rời khỏi đây thôi,'' Reiko-san đột ngột lên tiếng, giọng khàn đặc. ``Nơi này làm mình cảm thấy buồn nôn\dots''

``Cậu ấy nói đúng,'' em tôi gật đầu, giọng run run. ``Nếu có thêm một đợt sạt lở nữa\dots''

Cô không nói hết câu, nhưng tất cả đều hiểu. Chúng tôi không thể để bất cứ ai trong chúng tôi gặp nguy hiểm nữa.

Chúng tôi vội vã di chuyển, cố gắng tránh những đống đổ nát và tìm kiếm một lối thoát. Mỗi bước đi đều khiến tim tôi thắt lại - vừa lo lắng cho hai người bạn còn mất tích, vừa sợ hãi trước khả năng tòa nhà có thể sập bất cứ lúc nào.

Trước mặt chúng tôi, một cánh cửa trượt bằng gỗ hiện ra trong ánh sáng mờ ảo. Nó giống hệt những cánh cửa quen thuộc ở hành lang trường Aogami, nhưng phía sau nó chỉ là một màn đen dày đặc, không thể nhìn thấu qua được.

\pov{Haragen Saikyou}

``Để đó, em mở cho.''

Tôi đưa tay kéo cánh cửa gỗ trượt trước mặt, nhưng nó chẳng hề nhúc nhích. Tôi nghiến răng, thử thêm lần nữa, trượt mạnh hết sức, nhưng vẫn im lìm, như thể có thứ gì ghì chặt từ bên trong.

``Chó đẻ thật đấy\dots'' tôi thốt lên bực dọc trong tiếng rít.

Reiko-san dè dặt: ``Cậu\dots\ có chắc là nên mở không?''

Tôi hít một hơi, rồi buông tay ra, đáp gọn: ``Cậu nghĩ chúng ta còn đường khác có thể đi được không?''

Nii-chan và cô bạn đồng thời quay đầu nhìn lại con đường đã đi qua. Bóng tối dày đặc như một bức tường đen ngòm chắn ngang, nuốt chửng mọi thứ phía sau lưng chúng tôi. Quay đầu lại lúc này đồng nghĩa với việc có thể sẽ lạc vào một mê cung không có điểm dừng, và có lẽ sẽ không bao giờ tìm được lối ra.

Không cần chờ hai người cất lời, tôi cũng đã biết rõ câu trả lời.

``Vậy thì\dots\'' tôi hít một hơi, dần lùi lại. ``Không mở được thì phá thôi.''

``Saikyou à, anh không nghĩ--'' Nii-chan khẽ gọi, như muốn can ngăn, nhưng tôi đã lấy đà, tung chân đạp thẳng.

\textbf{*RẦM!*}

Tiếng gỗ gãy răng rắc vang lên chát chúa, cánh cửa bật ngược ra, đổ sập xuống nền.

Cửa bật tung, trượt ra ngoài ray, ngã sập xuống sàn gỗ trong tiếng va đập khô khốc. Reiko-san giật mình kêu khẽ, còn Nii-chan thì cau mày, bất ngờ nhìn tôi: ``Em lại dùng cách đó nữa hả?''

Tôi phì cười, nhún vai: ``Nii-chan, có cách nào nhanh hơn đâu.''

Cánh cửa giờ nằm chỏng chơ, mở ra một hành lang dài hun hút. Ánh sáng đỏ rực từ dãy đèn huỳnh quang treo trần hắt xuống, khiến không gian loang lổ sắc máu. Tôi nheo mắt: các bóng đèn nhấp nháy bất thường, có cái nhòe mờ, có cái lại chớp tắt liên hồi, như đang hấp hối. Tường, trần và sàn loang lổ chi chít những vệt đen như vết cháy, kéo dài ngoằn ngoèo khắp mọi nơi, trông như nội tạng đang bị nhuộm đỏ, nhưng đồng thời lại giống như thứ gì đang thấm ra từ bên trong. Cảnh tượng ấy khiến tôi thoáng rùng mình, cảm giác chẳng khác nào bước thẳng vào một cơn ác mộng sống sờ sờ.

\img{e1-8g}

Reiko-san khẽ rùng mình, tay bấu chặt đèn pin: ``Nơi này\dots\ trông như\dots\ một bể máu vậy\dots''

Tôi gật nhẹ, cố giữ giọng bình thản: ``Cẩn thận. Đi sát nhau thôi.''

Chúng tôi bước vào. Mỗi bước chân vang lên, dội lại trong không gian đỏ máu. Tôi liếc sang Nii-chan, thấy anh ấy im lặng, nhưng ánh mắt không giấu được sự cảnh giác. Cô bạn tóc hai bím thì thở dồn dập, chỏm tóc ngố trên đỉnh run rẩy, đôi lần còn quay lại nhìn phía sau, dù biết chẳng có gì ngoài bóng tối.

Đi được một đoạn, Nii-chan lẩm bẩm: ``Cảm giác như\dots\ chúng ta đã đi qua nơi này rồi.''

Tôi nhìn lên trần, nơi những đốm đèn đỏ chớp tắt. Trong khoảnh khắc ấy, ký ức tràn về, đúng, tôi và anh trai từng thấy khung cảnh này trong giấc mơ, và thứ sắc đỏ này tôi đã thấy trong phòng Mỹ thuật. Tôi nuốt khan, không dám nói rõ.

Khi cả ba tiến tới giữa hành lang, Reiko đột ngột dừng bước, ánh đèn pin của cô ấy run rẩy lia về phía trước. Tôi và Nii-chan lập tức nhìn theo.

Trước mặt, chắn ngang gần hết lối đi, là một đống bàn ghế học sinh. Chúng được chồng chất lên nhau một cách lộn xộn, tạo thành bức tường ngổn ngang, những chiếc chân ghế cắm chĩa ra như móng vuốt.

Tôi lạnh sống lưng. Cảnh tượng này\dots\ y hệt cơn ác mộng đã ám lấy tôi và Nii-chan. ``\dots Thật luôn?''

Reiko-san thốt lên, giọng không khỏi khó hiểu xen lẫn sợ sệt: ``Không thể nào\dots\ tại sao\dots\ lại có thứ này ở đây?''

``Tôi cũng chẳng biết,'' tôi thở dài mệt mỏi. ``Nhưng mà nếu chúng ở đây thì có lẽ phải nhấc chúng ra thôi.''

Tôi hít một hơi, định bước tới nhấc thử đống ghế kia xem có di chuyển được không. Mới đi được ba bước thì *RẦM!*, cánh cửa lớp học ngay bên cạnh bất ngờ trượt xuống, va mạnh xuống sàn, khiến tôi giật nảy, tim như nhảy khỏi lồng ngực.

``Eeek!''

``Cái quái\dots!?''

Reiko-san hét khẽ, lùi ngay lại, còn Nii-chan lập tức dang tay chắn trước, mắt đảo quanh. Tôi cũng sững sờ vài giây, rồi siết nắm tay, nuốt khan, cố giữ bình tĩnh.

``Có lẽ\dots\ chỉ là cánh cửa cũ quá thôi. Bình tĩnh nào, bình tĩnh\dots'' tôi lên tiếng trấn an.

Chúng tôi từ từ tiến lại phía lớp học vừa đổ cửa. Chưa kịp nhìn rõ bên trong thì *VÚT!* một chiếc ghế từ trong lớp bất ngờ bị ném ra với tốc độ rất nhanh, găm thẳng vào cửa lớp đối diện, chỉ cách đầu Reiko-san vài phân. Gió rít sượt qua, để lại tiếng va chát chúa làm cả ba chết lặng.

``\textbf{Gyaaaah!!!}'' *BỤP*

Cô ấy sụp thẳng xuống sàn, hai tay ôm chặt hạ bộ, mặt trắng bệch như giấy. Rõ ràng cậu ấy vừa trải qua khoảnh khắc ``suýt nữa thì\dots'' và đang phải gồng hết sức để giữ lại.

``\textbf{Reiko-san!}'' cả hai anh em đồng thanh gọi. ``Cậu ổn chứ!?''

``M-Mình\dots\ Mình vẫn ổn\dots'' cô ấy lắp bắp, mắt rơm rớm nước. ``X-Xém nữa thì\dots''

Nii-chan vội đỡ cô đứng dậy, thì thầm trấn an, may mà chưa có ``tai nạn'' nào khác xảy ra. Tôi nhìn cảnh đó mà vừa thương vừa thở phào, rồi gật nhẹ với anh.

``Anh lo cho Reiko-san đi. Nii-chan, đưa đèn pin cho em.''

``Ừm.'' Anh ấy gật, trao chiếc đèn cho tôi, rón rén đi về trước. Tôi siết chặt tay cầm, ánh sáng yếu ớt rọi vào khoảng không đỏ rực phía trước. Tôi dẫn đầu, bước vào lớp học, trong khi Reiko-san còn run rẩy nép vào cạnh anh tôi.

Bên trong, sắc đỏ tiếp tục nuốt chửng cả không gian. Các dãy bàn ghế chất đống ngổn ngang, một vài chiếc lật ngửa, chân ghế chĩa loạn xạ như bẫy sắt. Trên tường loang lổ vô số vệt đen, có cái chảy dài như vết nhầy, có cái như thể bị đốt cháy.

Ở giữa lớp, một bóng người ngồi bất động trên ghế. Tóc gáy tôi dựng đứng. Tôi lia đèn pin thẳng vào, nhưng bóng đó vẫn không cử động.

Chúng tôi vừa bước gần thêm vài bước, thì nó bỗng giãn ra, cơ thể vặn vẹo, tay chân uốn éo một cách méo mó, như hình ảnh bị lỗi. Tôi hét khẽ, còn Reiko-san thì hoảng sợ, bám chặt lấy cánh tay Nii-chan, rúc đầu vào ngực anh ấy để tránh nhìn cảnh tượng kinh hoàng trước mắt. Nii-chan siết vai tôi, cả ba đều lùi lại bản năng, không ai dám ho he lấy một câu.

Rồi trong chớp mắt, bóng người ấy biến mất, tan vào không khí.

Trên chiếc ghế nơi nó vừa ngồi, có một vật màu xanh lá, to cỡ nửa gang tay, nổi bật giữa nền đỏ. Reiko-san run rẩy chỉ vào: ``Cái\dots\ cái đó là\dots''

Tôi tiến lại gần, nhặt lên. Trông nó giống như ột bảng mạch, mặt kim loại sáng nhạt, chi chít linh kiện nhỏ. Không có pin như thứ trắng bạc tôi từng thấy trước đó, nhưng rõ ràng là một phần quan trọng.

``Nii-chan, nhìn này,'' tôi chìa cho anh.

Anh lật mặt sau, đọc mẩu chú nhỏ khắc bằng số: ``Có số sê-ri\dots\ và\dots\ dường như khớp với phần vỏ sau của chiếc đồng hồ đeo tay. Saikyou, em có thể ướm được không?''

Tôi lập tức lôi phần vỏ ra khỏi túi váy, áp thử. Quả thật, đường viền khớp hoàn hảo. Một cảm giác quen thuộc dấy lên trong lòng tôi, sự thôi thúc kỳ lạ mà mỗi khi tìm được bộ phận của chiếc đồng hồ đeo tay, tôi không thể nào phớt lờ.

``Để em cầm nó cho,'' tôi nói chắc nịch, trượt bo mạch vào túi cùng vỏ sau.

Reiko-san vẫn còn run rẩy, nhưng ánh mắt cô ấy sáng lên đôi chút, như vừa tìm thấy tia hy vọng giữa ác mộng. Anh trai tôi gật nhẹ, không có ý phản đối.

Như vậy, chiếc đồng hồ đeo tay kia đang dần dần thành hình - hai dây đeo, hai vỏ bọc, đế sạc, dây cáp và bo mạch chủ đã có đủ, dù vẫn thiếu những bộ phận cốt lõi để có thể hoạt động.

Không chần chừ thêm, chúng tôi rời khỏi lớp học, quay lại hành lang đỏ rực đầy vệt đen. Bóng tối và tiếng ù ù vẫn đè nặng, càng thôi thúc chúng tôi đi nhanh hơn.

\ct

Chúng tôi tiếp tục bước đi, bám sát theo ánh đèn pin trong tay Nii-chan. Không gian trước mắt vẫn đỏ đặc quánh, tường, sàn, trần đều phủ chằng chịt những vệt đen ngoằn ngoèo kéo dài vô tận. Thỉnh thoảng, tiếng lách tách từ những bóng đèn huỳnh quang trên cao vang lên, khiến tôi rùng mình, cứ như thể nơi này đang cố thoi thóp thở.

Reiko-san đi sát ngay cạnh tôi, mắt vẫn đỏ hoe, nhưng cậu ấy cắn môi, cố giữ cho đôi chân không run rẩy. Tôi không dám buông tay khỏi tay cậu ấy, vì nếu không, có lẽ cậu ấy sẽ không thể đi được. Trong khi đó, Nii-chan đi đầu, từng bước thận trọng, vừa rọi đèn vừa ngoái lại nhìn chúng tôi.

``N-Này hai người\dots'' cô bạn tóc hai bím chợt lên tiếng.

Cả hai chúng tôi ngước lên phía trước. Chúng tôi nhận thấy con đường phía trước sạch sẽ hơn, những vệt đen ghê rợn kia biến mất hẳn, thay vào đó là nền tường và sàn bình thường, tuy chìm trong bóng tối.

Tôi lập tức chỉ tay: ``Nhìn kìa, hết rồi\dots''

Reiko-san cũng thở ra một hơi thật nhẹ, đôi vai dường như giãn ra sau bao nhiêu căng thẳng. Nhưng cả ba chúng tôi đều biết, bình thường không có nghĩa là an toàn.

Chúng tôi men theo hành lang mới này, bước chậm rãi trong thứ bóng tối đặc quánh, chỉ nghe tiếng bước chân mình dội lên từng nhịp. Không gian im ắng đến nghẹt thở.

Rồi ánh sáng bất chợt xuất hiện. Ở cuối dãy, một bóng đèn huỳnh quang trắng đang chập chờn, lóe tắt liên tục, hắt xuống một cánh cửa trượt gỗ. Ánh sáng mờ run rẩy khiến cánh cửa ấy nổi bật một cách kỳ dị, như đang chờ đợi ai đó có thể mở ra.

``Để anh mở.''

Nii-chan hít một hơi, chậm rãi đi lên trước. Bàn tay anh đặt lên tay nắm, đẩy ra. Không hề có sự kháng cự nào, cánh cửa mở trơn tru đến lạ, khác hẳn cánh cửa mà tôi đã mạnh chân đạp không thương tiếc.

Bên kia là một hành lang khác, sáng choang bởi những bóng huỳnh quang trắng tinh. Tường trắng, sàn sạch, không có lấy một vệt đỏ hay đen nào. Mọi thứ trông\dots\ bình thường một cách khó tin.

Nhưng tôi không thấy nhẹ nhõm. Việc cánh cửa ấy mở ra một cách dễ dàng như vậy, và mọi thứ bình thường như thế\dots\ hẳn là có điềm gì đó chẳng lành. Chúng còn khiến chúng tôi bất an còn hơn cả hành lang đỏ máu ban nãy.

``\dots Đi thôi nào,'' tôi nhẹ giọng. Cả hai người gật đầu, tiếp bước theo sau.

Vẫn như ban nãy, cả ba chúng tôi bước đi, nhịp chân chậm mà đều, từng tiếng vang vọng dội lại giữa hành lang trắng sáng. Không còn mùi sắt gỉ hay những vệt đen vô tận nữa, chỉ là ánh sáng huỳnh quang, tường trơn nhẵn, sạch sẽ đến mức đáng ngờ. Càng đi, tôi càng thấy như có thứ gì đó treo lơ lửng trong không khí, ép chúng tôi phải nói thật khẽ, thật dè chừng.

``Cứ như thể\dots\ ai đó đang dõi theo mình vậy,'' Reiko-san thì thầm, mắt liếc về sau.

Nii-chan cười nhẹ, nhưng nụ cười không hề tự nhiên. ``Nếu có, thì chắc hẳn họ đã thấy rõ chúng ta đang làm gì. Em nghĩ sao, Saikyou?''

Tôi khẽ gật, mắt vẫn đảo qua từng bóng đèn, từng khe nứt nhỏ trên tường. ``Cảm giác này\dots\ không giống ảo giác đâu. Có gì đó sắp tới. Một cái gì đó kinh khủng.''

Những bước chân tiếp tục dẫn chúng tôi đi qua đoạn hành lang dài, ánh sáng trắng lặng lẽ phủ lên cả ba như một lớp màn mỏng. Nó chẳng khác một hành lang trường học về đêm: trống trải, yên tĩnh, nhưng ẩn chứa một thứ gì đó muốn bật ra bất cứ lúc nào.

Và ngay khi tôi đã xác định rằng điều \emph{tồi tệ} đó sẽ xảy ra\dots

*Ù Ù Ù\dots*

Bất ngờ, mặt đất rung lên. Một cơn địa chấn dữ dội ập đến, khiến tôi khụy gối xuống theo bản năng. Những mảng trần rung lắc, đèn huỳnh quang chao đảo như sắp rơi. Toàn bộ hành lang lắc lư, như thể cả tòa nhà này đang bị gió cuốn đi.

``\textbf{Giữ chặt vào!}'' tôi hét, hai tay bám vào tường, cảm nhận từng đợt chấn động truyền qua cơ thể.

Reiko-san thở gấp, cố nắm lấy vai tôi. ```Chuyện gì vậy!? Động đất\dots\ ở đây sao!?''

Nii-chan nghiến răng, đôi mắt căng ra, bám víu lấy hai người chúng tôi. ``Khỉ thật chứ\dots\ Đến cả nơi này cũng có sao!?''

Trận động đất mạnh đó chỉ chỉ kéo dài vài giây, rồi mặt đất ngừng hẳn. Nhưng với chúng tôi, nó dài như cả một thế kỷ. Tiếng kêu rên của bê tông và sắt thép lắng xuống, chỉ còn lại ánh đèn huỳnh quang chập chờn, lúc sáng lúc tối. Cả ba chúng tôi dồn lại gần nhau, mồ hôi lạnh rịn ra nơi thái dương.

Và rồi, những bóng dáng đen xuất hiện. Vài cái bóng người hốt hoảng chạy ngang qua hành lang phía trước, ánh mắt như chỉ hướng tới lối thoát vô hình nào đó. Không ai quay lại, chỉ có những bước chân vội vã, những tiếng hét thất thanh nghe như của con gái dần biến mất vào xa.

``\textbf{Đợi đã!}'' tôi hét lên. ``\textbf{Có người kìa!}''

Ngay lúc đó, chiếc bóng đèn huỳnh quang trên trần nhà trước mặt chúng tôi bỗng rung lên bần bật. Nó lắc lư dữ dội trong vài giây, rồi đứt phựt khỏi giá treo, rơi thẳng xuống sàn với một tiếng *XOẢNG!* chói tai. Mảnh thủy tinh vỡ văng tung tóe, ánh sáng chập chờn rồi tắt ngấm.

``Cẩn thận!'' Kyuen-san kéo tôi lùi lại, tránh những mảnh vỡ sắc nhọn.

Như một bản năng sinh tồn, chúng tôi lao lên, bám theo quỹ đạo của những bóng dáng ấy, cho tới khi hành lang chia ra một ngã ba. Bên phải là một cánh cửa trượt đóng kín, bên trái là hành lang sáng đèn kéo dài vô tận, theo hướng mà các bóng dáng đen ấy vừa chạy. Cả ba chúng tôi khựng lại, ánh nhìn trao đổi trong khoảnh khắc căng thẳng tột độ.

Rồi, như một sự đồng thuận thầm lặng, chúng tôi rẽ sang bên trái. Hành lang trước mặt sáng đèn đều đặn, nhưng không ai trong chúng tôi thấy yên tâm.

``Mình nghĩ\dots\ chắc họ chạy tới cánh cửa kia,'' Reiko-san chỉ về phía trước.

Chúng tôi mới chỉ bước được vài nhịp thì đèn vụt tắt. Từng bóng một, từ cuối hành lang phía xa dần dần tắt ngúm, tiến lại gần chỗ chúng tôi đứng. Bóng tối nuốt chửng ánh sáng, từng đoạn, từng đoạn, cho đến khi hành lang rơi vào một khoảng im lặng đặc quánh.

``C-Cái gì vậy?'' cô nàng mắt hai màu rít lên, giọng run rẩy. ``Sao đèn lại tắt hết thế này?''

``Đừng hoảng,'' Nii-chan nói, cố trấn an mặc dù giọng cậu cũng không được vững. ``Chắc chỉ là mất điện thôi.''

Tim tôi đập mạnh. Chúng tôi nín thở. Và rồi, *CẠCH!*, một tiếng động vang lên, vọng từ phía hành lang đã chìm trong tối.

``\textbf{Có ai ở đó không?}'' tôi cất tiếng hỏi, cố không để nỗi sợ lộ ra trong giọng nói.

``Để mình rọi đèn xem,'' Reiko nói, giọng căng thẳng. Cô hướng chiếc đèn pin rọi về hướng đó, và\dots\ suýt buông thõng cả tay.

Một cô gái (hoặc một thứ gì giống cô gái) vụt qua ánh sáng. Gương mặt cô ta\dots\ không còn nguyên vẹn, xộc xệch như thể bị ai đó xé toạc ra, nhưng đôi mắt vẫn mở trừng trừng. Cơ thể mờ ảo, gần như trong suốt, chạy sượt qua ngay trước mắt chúng tôi, không một lời, không một âm thanh.

``\textbf{Cái\dots\ cái gì vậy!?}'' Reiko hét khẽ, tay bấu chặt lấy vai anh tôi.

``Bình tĩnh nào, chỉ là ảo giác thôi!'' Nii-chan nghiến răng, mặc dù tôi có thể thấy anh ấy cũng đang run.

Còn bản thân tôi\dots\ tôi chỉ có thể lẩm bẩm: ``Chuyện gì vừa xảy ra vậy\dots?''

Dẫu miệng nói ra thế, lòng tôi biết rõ: đây không phải ảo giác. Và tim tôi vừa bị hù cho một nhát chí mạng.

Chúng tôi quyết định tiến lên phía trước. Nhưng càng đi, cánh cửa trượt cuối hành lang chẳng hề gần lại. Cứ như thể hành lang này cố tình kéo dài ra, nuốt chửng từng bước chân, khiến khoảng cách kia mãi chỉ là một ảo vọng xa vời. Tôi có thể cảm nhận được mồ hôi lạnh chảy dọc sống lưng, từng hơi thở nặng nề dần.

``Kỳ lạ thật,'' anh tôi lẩm bẩm, giọng khàn đặc. ``Rõ ràng chúng ta đã đi được một đoạn dài, nhưng sao cánh cửa vẫn xa như vậy?''

``Em cũng thấy thế,'' tôi vừa bước đi, vừa cảm thấy mệt nhọc tói lạ. ``Cứ như là\dots\ hành lang này đang\dots\ đùa với chúng ta vậy.''

``Không thể nào,'' Reiko cố trấn tĩnh. ``Chắc chắn là do chúng ta mệt mỏi quá thôi. Cố lên một chút nữa.''

\dots

\dots

\dots

Sau một hồi đi vô ích, cả ba đành dừng lại, thở dốc trong tuyệt vọng.

``Vô ích rồi,'' Nii-chn thở hổn hển. ``Chúng ta phải quay lại thôi.''

``Anh nói đúng,'' Saikyou gật đầu, mắt nhìn về phía sau. ``Em có cảm giác\dots\ chúng ta nên đi ngược lại.''

Không còn lựa chọn nào khác, chúng tôi quay đầu, quay về hướng ban đầu.

Và bất ngờ thay, ở trước mặt, một cánh cửa trượt khác, mở hé, như đang chờ sẵn.

\pov{Haragen Reiko}

Cả ba khựng lại, ánh nhìn trao đổi trong khoảnh khắc căng thẳng tột độ. Qua khe hở nhỏ của cánh cửa, một luồng gió lạnh lẽo thoát ra, mang theo mùi ẩm mốc kỳ lạ khiến tóc gáy tôi dựng đứng. Tôi quyết định bước lên trước.

``Để\dots\ để mình mở cho,'' tay tôi khẽ chạm vào cánh cửa trượt. Nó kêu két một tiếng nhỏ rồi tách ra, để lộ một căn phòng nhỏ.

Bên trong giống như một văn phòng bị bỏ quên. Ghế chồng chất lên nhau, nằm ngổn ngang, ngả nghiêng như sắp đổ. Một khung toan trắng tinh đặt ngay chính diện, bên cạnh là hộp dụng cụ vẽ mở hé, màu sắc lấm lem. Không khí trong căn phòng lạnh hơn hẳn ngoài hành lang, và im lặng đến đáng sợ.

Tôi quay đầu lại, trên mép cửa có dán một mảnh giấy. Những dòng chữ nguệch ngoạc hiện ra:

\txtbx{Đầu mình đau quá\dots\\
  Chuyện gì đang xảy ra với mình vậy\dots?\\
  Mình cứ vẽ, cứ vẽ mải miết như thế\dots\ Nhưng tại sao nó lại không ra được tác phẩm như mong muốn vậy?\\
  Không lẽ\dots\ trình độ của mình chỉ có thế?}

Một thoáng rùng mình chạy dọc sống lưng. Những lời này nghe thật tuyệt vọng, như lời trăng trối. Tôi đưa mắt nhìn Kyuen và Saikyou; cả hai đều cau mày, như đã linh cảm được điều gì bất thường.

``Những dòng chữ này\dots\ có phải là của cô ấy không?'' Saikyou-san nghiêng đầu, nhìn tờ giấy.

``Cô gái vẽ tranh mà mọi người vẫn đồn thổi ấy à?'' Kyuen-san tiếp lời, mắt vẫn dán chặt vào mảnh giấy, rồi nhắm mắt lại như đang cố nhớ lại. ``Người ta nói cô ấy đã\dots\ dùng chính máu của mình để hoàn thành kiệt tác cuối cùng.''

``Có lẽ\dots\ là cô ấy đấy,'' tôi gật đầu, giọng trầm xuống. ``Những lời này\dots\ như thể cô ấy đang dần mất đi lý trí. Sự ám ảnh với nghệ thuật đã đẩy cô ấy đến bờ vực của điên loạn.''

Tôi cảm thấy không khí xung quanh như lạnh đi vài độ. Câu chuyện về một nữ sinh mỹ thuật tài năng nhưng đã cống hiến bản ``kiệt tác'' của mình theo hướng cực đoan nhất - những lời đồn ấy giờ dường như đang hiện hữu ngay trước mắt chúng tôi.

Và rồi\dots\ nó xảy ra.

Bỗng nhiên, một cơn nhói nhẹ lan khắp đầu tppi. Không rõ từ đâu, nhưng nó khiến tôi khẽ nhăn mặt, một tay đưa lên thái dương. Chưa kịp cất tiếng, Saikyou-san đã khụy xuống, hai tay ôm chặt lấy đầu, hơi thở dồn dập.

``Nii-chan\dots\ sao tự dưng đầu em đau quá vậy\dots!'' giọng cậu run rẩy, nhưng mắt vẫn dán chặt vào khung toan trắng.

Tôi liếc xuống khung vẽ ấy, và rồi nhận ra\dots\ nó đã không còn trắng nữa.

Chúng tôi đứng sững trước một bức họa đầy ám ảnh - những nét chì đen đan xen với những vệt màu đỏ tươi như máu đang rỏ giọt. Trong tranh là hình ảnh một cô gái trẻ, ánh mắt mở trừng trừng nhìn thẳng về phía chúng tôi. Trên môi cô ấy là một biểu cảm kỳ lạ, nửa như nụ cười méo mó, nửa như tiếng thét kinh hoàng bị đóng băng - không ai trong chúng tôi dám chắc đó là gì.

\dots Liệu đó có phải là bức tranh chết chóc\dots\ mà các bạn trong trường nói tới không?

Tim tôi đập loạn. Một luồng cảm xúc lạnh lẽo, nghẹt thở xuyên thẳng vào ngực. Cảm giác uất ức, tuyệt vọng, trĩu nặng như chì đè ép tâm can.

``\dots Không phải\dots\ \textbf{lại nữa chứ!!!}'' Saikyou-san hét lớn, giọng nghẹn lại, tay bấu chặt vào ngực áo. Tôi cũng cảm nhận được, nhưng ở em, nó mạnh mẽ gấp bội. Cảm giác như toàn bộ bóng tối của bức tranh đang nhấn chìm lấy cậu ấy.

``Không được! Ta phải rời khỏi đây thôi!'' người anh hét lớn, lao tới. Cậu nắm chặt tay tôi và em gái cậu, kéo phăng ra khỏi căn phòng.

Ngay khi bước chân ra ngoài, cơn đau đầu lập tức biến mất như chưa từng xảy ra. Ba đứa chúng tôi ngã quỵ xuống nền hành lang, thở hổn hển, mồ hôi chảy ròng ròng.

``Suýt nữa thì\dots'' tôi cố hít sâu, vẫn còn thấy ngực quặn thắt.

Tôi thấy Saikyou-san rúc đầu vào ngực anh trai, vai run lên từng đợt. Kyuen-san vòng tay ôm lấy em gái, ánh mắt đau đớn nhìn về phía trước. Nhìn cảnh tượng ấy, tôi chợt hiểu được phần nào vì sao cô em gái của cậu lại ám ảnh với bức tranh đó đến vậy. Đến bản thân tôi cũng cảm nhận được rằng bức tranh ấy đã không bình thường, như thể nó là một lời cảnh báo cho những gì sắp xảy ra.

Cảnh tượng trước mắt tôi\dots\ vừa choáng ngợp, vừa khiến lòng bàn tay tôi lạnh toát. Và kỳ lạ thay, trong sâu thẳm, tôi cảm thấy nó\dots\ rất quen thuộc. Giống như thể một người đã tìm thấy bến đỗ mới trong đời mình vậy.

Tôi ngước về phía trước, rồi khẽ giật mình với quang cảnh vừa thực vừa ảo. Tôi khẽ gọi hai người ấy: ``Kyuen-san, Saikyou-san\dots\ nhìn kìa.''

Họ cũng đều khựng lại, ánh nhìn trao đổi trong khoảnh khắc ngỡ ngàng.

Một con đường dài hiện ra, sáng lên trong bóng tối mịt mùng. Hai bên là những bụi hoa bỉ ngạn đỏ rực như máu, gió thổi qua khiến từng cánh hoa khẽ lay động. Con đường trông không giống như được lát đá, mà phản chiếu như mặt nước, lấp lánh những gợn sóng nhỏ. Cuối con đường ấy, một cổng torii đỏ sẫm vươn cao, sau đó là một ngôi đền, được ánh sáng từ trên cao chiếu xuống, khiến nó trở nên vừa thiêng liêng, vừa đáng sợ.

\img{e1-8h1}

Tôi khẽ thì thầm, ``\dots Chẳng phải chỗ này\dots''

Cậu con trai cau mày: ``\dots là đường mà Sora-san đi lên đền sao?''

Em gái của cậu ôm chặt cánh tay mình, tiếp lời: ``Không thể nhầm vào đâu được. Hình như\dots\ chúng ta đã từng đi con đường này.''

Chúng tôi bước chậm rãi về phía ngôi đền, lòng bàn chân như dẫm trên thứ gì đó mềm mại, nhưng khi cúi xuống, vẫn chỉ thấy mặt đường loang loáng phản chiếu. Mỗi bước đi, âm thanh vang vọng như thể có hàng trăm người khác cũng đang bước cùng.

Toàn thân tôi như nặng trĩu. Một cảm giác quen thuộc đến mức đau đớn, đúng vậy, nơi này chính là nơi Sora-san từng tới để ước nguyện. Giờ đây, bọn tôi đang lặp lại con đường đó, mà không rõ vì sao.

Khi tiến gần đến ngôi đền, ánh sáng hắt rõ bóng một người đang đứng quay lưng về phía chúng tôi. Mái tóc nâu dài thả nhẹ, trên đầu cài một chiếc kẹp tóc đỏ lệch sang trái, khẽ rung theo làn gió.

\img{e1-8h2}

Tôi sững lại, đôi môi run run, chỉ tay về phía trước: ``Kia\dots\ chẳng phải là\dots''

Cô gái ấy khẽ xoay người lại. Đôi mắt xanh nước biển sáng lên trong bóng tối, sâu thẳm và xa xăm, như thể xuyên qua cả thực tại này.

Không thể nhầm lẫn được.

Misora Shino-san.

\pov{-}

Reiko khựng lại, đôi mắt cô mở to khi nhìn thấy dáng hình quen thuộc đang đứng trước mặt. Cổ họng ứ lại, nhưng khóe môi cô thoáng nở một nụ cười run rẩy.

``Shino-san\dots?'' từng tiếng cô bật ra như một hơi thở nhẹ. Trong thoáng chốc, mọi gánh nặng trong lòng như tan biến, chỉ còn sự nhẹ nhõm vì cuối cùng cũng gặp lại người bạn thân của mình.

Thế nhưng, người bạn trước mặt lại có gì đó khác lạ. Đôi mắt xanh trong trẻo ngày nào giờ phủ một tầng sương buồn, ánh nhìn khiến Reiko chùng xuống. Kyuen và Saikyou cũng nhận ra điều đó - không còn cảm giác sát khí lạnh lẽo như lần đầu gặp mặt, mà thay vào đó là một nỗi u buồn sâu thẳm, như thể cô gái trước mặt họ đang mang trên vai cả một bầu trời đau thương không thể nói thành lời.

``Cậu\dots\ có thực sự ổn không vậy?'' cô dè dặt, bàn tay hơi đưa về phía trước, như sợ chạm vào một giấc mơ mong manh.

Shino cúi mặt, môi run nhẹ. Nụ cười cô hé ra nhưng không còn rực rỡ như trước. ``Reiko-san\dots\ Mình không ổn đâu. Bởi vì mình\dots\ đã chết rồi.''

Lời đáp buốt lạnh như gió xuyên thẳng vào ngực Reiko. Cô tròn mắt, lùi lại nửa bước.

``Cậu nói\dots\ gì cơ? Không thể nào\dots\ Mình tưởng\dots\ cậu chỉ là\dots\ bị lạc\dots\ chỉ tự nhận là Sora để trốn tránh thôi chứ\dots?''

Shino khẽ lắc đầu. Đôi mắt xanh dường như còn sáng hơn khi đẫm buồn. ``Ước gì mọi chuyện chỉ đơn giản vậy. Nhưng không\dots\ trận sạt lở đất hôm đó, mình đã thật sự bỏ mạng.''

Reiko choáng váng, bàn tay siết chặt ngực áo.  Tới lúc này, cô mới thực sự tin cô bạn ấy đã không còn trên đời nữa. Không khí như đông cứng lại, cho tới khi Kyuen và Saikyou bước tới, sắc mặt nghiêm trọng.

Cậu con trai lên tiếng, giọng trầm hẳn xuống: ``Chúng tôi\dots\ đã thực sự biết được chuyện gì đã xảy ra lúc đó.''

Cô em gái của cậu gật đầu, chậm rãi nói thêm: ``Cậu nói đúng. Đó thực sự là một trải nghiệm kinh hoàng. Đến cả hai anh em tôi suýt bỏ mạng ở đó.''

Nghe vậy, Shino khẽ nghiêng đầu, ánh mắt như đã đoán được từ trước. Cô thở dài, rồi nhìn lại Reiko vẫn đang run rẩy, chưa thể chấp nhận được sự thật phũ phàng ấy.

``Các cậu thấy đấy\dots\ chuyện đó không chỉ xảy ra một lần. Nó đã lặp lại\dots\ hết lần này đến lần khác. Và mỗi lần\dots\ đều kết thúc bằng cái chết của tất cả. Như một lời nguyền quái ác vậy.''

Kyuen siết chặt nắm tay, giọng đanh lại: ``Tại sao\dots\ lại như thế?''

Cô nàng mắt xanh ngẩng lên, đôi mắt ánh lên sự mỏi mệt. ``Nó bắt nguồn từ sự kiện kinh hoàng hôm ấy, như các cậu đã từng chứng kiến. Và\dots\ nó gắn liền với học sinh chuyển đến.''

Saikyou nhíu mày, bước tới một nửa, giọng thấp xuống: ``Ý cậu là sao? Học sinh chuyển đến\dots\ có liên quan gì?''

Shino khẽ cười, một nụ cười pha lẫn chua xót. ``Các cậu cũng biết rồi\dots\ mình luôn khao khát có một người bạn, phải không? Thế nhưng\dots\ mong ước ấy đã dần bị bóp méo. Thành một vòng lặp. Mỗi khi có người mới đến\dots\ bi kịch ấy lại tái diễn, y như cũ. Không ai thoát được.''

Lời cô buông xuống như một hồi chuông tang, kéo tất cả rơi vào im lặng. Reiko chỉ còn biết siết chặt đôi bàn tay run rẩy, lòng quặn thắt khi thấy bạn mình vừa gần gũi, vừa xa vời không thể chạm tới.

``Shino-san\dots\ tại sao lại thành ra thế này\dots'' tiếng cô nghẹn lại, như sắp bật khóc.

Cô nàng tóc nâu vòng tay ôm lấy Reiko, dịu dàng nhưng ánh mắt vẫn phủ một màn u uất. ``Xin lỗi cậu, Reiko-san. Lẽ ra mình phải giải thích từ sớm\dots''

Cậu con trai khoanh tay lại, ánh nhìn nghiêm nghị xoáy vào Shino. ``Nhưng mà tôi thắc mắc\dots\ thứ đang kéo chúng tôi vào vòng lặp này\dots\ là cái gì?''

Cô nàng mắt xanh im lặng một thoáng, rồi cất giọng đều đều, như hỏi lại để xác nhận: ``Các cậu biết có một người trông y hệt mình, chỉ khác là đôi mắt nâu không?''

Saikyou chau mày: ``Cái người tự nhận là `Sora' sao?''

Cô nàng tóc nâu khẽ gật đầu. ``Đúng vậy. Đó chính là Sora. Nhưng thực thể đó không phải là một con người bình thường\dots\ mà là một phần ký ức bị tách khỏi mình.''

Reiko sững sờ, ngẩng lên nhìn cô bạn thân. ``Ký ức\dots\ bị tách ra?''

``Ừ.'' cô đáp, rồi cô trầm ngâm hẳn. ``Sora là hiện thân cho ước vọng kết bạn mãnh liệt của mình. Nhưng đồng thời\dots\ cũng là cách mình chối bỏ hiện thực. Mỗi khi có học sinh mới chuyển đến, vòng lặp này lại khởi động. Và mọi thứ\dots\ kết thúc bằng thảm kịch.''

Kyuen cau chặt mày, giọng đanh lại: ``Vậy tức là\dots\ chừng nào Sora còn tồn tại, vòng lặp sẽ còn tiếp diễn?''

Shino nhìn sang cậu ánh mắt lóe lên sự kiệt quệ: ``Đúng thế. Chừng nào thực thể ấy còn ở ngoài kia, cô ấy sẽ cứ lún sâu hơn vào nỗi đau. Không thể dứt ra. Mình đã nhiều lần cố kéo Sora quay lại, nhưng,'' cô khẽ lắc đầu, nụ cười chua chát thoáng hiện trên môi, ``\dots Sora quá cứng đầu. Cô ấy luôn nói rằng muốn `kết thân với các bạn', rằng đó là ước nguyện cuối cùng. Và cô ấy không chịu trở về.''

Saikyou chậm rãi lên tiếng, ngữ điệu pha lẫn bối rối và suy tư: ``Vậy\dots\ trong mắt bọn tôi, Sora chẳng khác nào một người thực sự tồn tại. Nhưng kỳ thực\dots\ chỉ là mảnh vỡ ký ức của cậu?''

Shino khép mắt lại, ôm Reiko chặt hơn như sợ mất đi chút hơi ấm cuối cùng. ``Không sai. Đó chỉ là một phần ký ức bị đánh cắp khỏi mình. Và càng kéo dài\dots\ mình càng mệt mỏi. Mình bắt đầu chán ngấy tất cả. Chán ngấy vòng lặp vô tận này.''

Trong vòng tay ấy, cô bạn cảm nhận rõ sự run rẩy mỏi mệt nơi người bạn mình. Trái tim cô nhói lên từng cơn, vừa thương vừa bất lực, không biết phải làm gì ngoài việc siết chặt cô bạn đã rời thế gian hơn nữa.

Shino khẽ thở dài, vòng tay ôm người bạn dần nới lỏng. Cô lùi lại nửa bước, ánh mắt xanh lặng lẽ lướt qua từng người một. Có gì đó trong thái độ và khí chất của cả ba khiến lòng cô rung lên một cảm giác mơ hồ, vừa xa lạ vừa có gì đó\dots\ tin cậy.

``Nhưng mà mình thấy có gì đó trong ba người các cậu\dots\ khác biệt hẳn so với những người mà mình đã từng gặp.''

Shino nhìn thật kỹ, rồi chậm rãi nói: ``Đầu tiên là Reiko-san. Có những khoảnh khắc mình cảm giác như cậu mang theo \emph{một sức mạnh nào đó để chống chọi lại vòng lặp này}, thứ mà mình chưa từng thấy, dù với cậu nó khá là mơ hồ.''

Reiko thoáng sững người, không biết phải trả lời thế nào. Cô chỉ mím môi, cúi đầu, còn bàn tay thì siết chặt lấy vạt áo, vì đến bản thân cô cũng không nhận ra là cô có làm vậy. Những lời Sakura nói từ trước bỗng chốc như văng vẳng bên tai, và có lẽ đó là điều mà cô ấy muốn nói tới.

Sau đó, ánh mắt của cô nàng dừng trên Saikyou. ``Còn cậu\dots\ Saikyou-san, phải không? Ở cậu có một sự điềm tĩnh lạ lùng. Cứ như thể cậu đã nhiều lần chứng kiến bi kịch này, và mỗi lần lại tích lũy thêm thứ gì đó\dots\ như \emph{một lớp ký ức chồng chất không thuộc về nơi này}.''

Cô gái tóc đuôi ngựa chau mày, định phản bác, nhưng lại thôi. Sự im lặng ấy chỉ khiến lời Shino càng thêm nặng nề.

Cuối cùng, cô xoay sang Kyuen. ``Và cậu\dots\ Kyuen-san. Ánh mắt cậu, nó khác hẳn với bất kỳ ai mình từng gặp. Cứ như thể cậu không hề sợ hãi trước cái chết, bởi vì nơi cậu đứng\dots\ \emph{vốn dĩ không phải ở đây}.''

Cậu con trai khẽ nghiêng đầu, nở một nụ cười nhạt: ``Nghe cậu nói chằng khác nào áp cho bọn tôi một cái mác lạc loài cả.''

Shino khẽ lắc đầu, giọng nhỏ nhẹ nhưng chắc nịch: ``Không phải mác bình thường đâu. Đó là sự thật. Cả ba người các cậu\dots\ giống như đến từ một nơi khác, vốn dĩ không thuộc về thế giới này. Chính vì vậy, vòng lặp vốn được định sẵn này đã bắt đầu trật khớp.''

Cả ba khựng lại, trao nhau những cái nhìn khó hiểu. Những lời mà cô ấy nói tuy rằng mơ hồ, nhưng không phải là không có cơ sở.

Reiko hít một hơi, lấy hết dũng khí để hỏi: ``Ý cậu là\dots\ chúng mình có thể thay đổi vòng lặp sao?''

Shino không đáp thẳng, chỉ khẽ nhếch môi, ánh mắt sâu thẳm. ``Có thể, mình cũng không chắc chắn. Nhưng nhìn những gì các cậu đã nhặt được những mảnh ghép dọc đường đi\dots\ mình tin là vậy.''

Saikyou cau mày: ``Những mảnh ghép\dots?''

``Những thứ giống như bộ phận của đồng hồ ấy,''  cô đáp, giọng cô như trôi dạt vào khoảng không, nửa như gợi ý, nửa như thử thách. ``Chúng không phải vật ngẫu nhiên. Chúng có thể giúp các cậu giữ nhịp cho chính mình, giữa một vòng lặp vốn chỉ muốn cuốn trôi tất cả.''

Kyuen nheo mắt, giọng mang chút nghi hoặc: ``Ý cậu là mấy thứ vụn vặt đó sẽ có vai trò như\dots\ trợ lý cho chúng tôi?''

Shino gật khẽ, nụ cười thoáng buồn: ``Đúng thế. Có thể chúng chính là chìa khóa. Nhưng chìa khóa để mở ra cánh cửa nào\dots\ thì các cậu sẽ phải tự mình tìm hiểu.''

Không gian trở lại tĩnh lặng, chỉ còn tiếng gió lùa qua cánh cổng torii. Trong ánh sáng mờ nhạt, lời nói của cô nàng tóc nâu như khắc sâu vào tâm trí cả ba người, một lời tiên tri vừa mơ hồ vừa ám ảnh, nhưng lại gieo rắc một niềm hy vọng mong manh.

Trong ánh sáng yếu ớt từ ngôi đền, gương mặt Shino khẽ dịu xuống. Đôi mắt xanh long lanh như phủ một tầng sương mờ, nhưng sau đó lại ánh lên một tia hy vọng mong manh.

``Các cậu\dots'' Shino ngập ngừng, rồi cuối cùng thốt ra, ``Nếu thực sự có ai có thể làm được\dots\ thì đó chính là ba người các cậu. Hãy thuyết phục Sora trở về với mình. Một mình mình đã không còn đủ sức nữa rồi.''

Đôi môi cô run run, ánh nhìn tha thiết như một lời nhờ vả, nhưng ẩn sâu trong đó lại là một niềm tin mãnh liệt, rằng ba người trước mặt thực sự có thể tạo ra sự khác biệt.

Reiko đứng lặng, những lời ấy đâm thẳng vào trái tim cô. Từ nỗi mông lung bấy lâu, một luồng nhiệt bừng cháy. Cô siết chặt nắm tay, ánh mắt sáng rực. ``Nếu có một chìa khóa\dots\ mình sẽ tìm ra. Và mình sẽ không bỏ cuộc. Đó có lẽ sẽ là di nguyện cuối cùng mà chúng mình sẽ thực hiện.''

Kyuen hít sâu, khóe miệng cong lên một nụ cười kiên nghị. ``Nghe thì giống như lao đầu vào địa ngục, nhưng được rồi. Nếu vòng lặp này thực sự có thể bị phá vỡ, tôi sẽ không ngồi yên nhìn nó nuốt chửng chúng ta thêm lần nào nữa.''

Saikyou khẽ gật đầu, giọng nhẹ nhưng dứt khoát:
``Em cũng thế. Nếu đó là cách để thoát khỏi bi kịch\dots\ thì em sẽ cùng hai người bước tiếp.''

Shino nhìn họ, nét mặt dần thả lỏng. Cô giơ tay, chỉ về phía cánh cổng của ngôi đền. Ngay lúc ấy, cánh cổng khẽ rung lên, rồi từ từ hé mở. Một luồng sáng trắng rực rỡ bắn thẳng vào màn đêm, phản chiếu xuống mặt nước, tạo thành con đường sáng dẫn lối.

``Khi bước qua đó\dots'' cô nói khẽ, giọng như vọng lại từ một nơi xa xăm, ``\dots cả ba người sẽ tới một thời điểm khác. Hãy chuẩn bị sẵn sàng.''

Ba người nhìn nhau, không ai lên tiếng, nhưng trong ánh mắt đều đã có câu trả lời. Quyết tâm hiện rõ: dù phía trước là gì, họ cũng sẽ cùng nhau đối mặt.

Khi cả ba bước tới gần cánh cổng, ánh sáng càng thêm chói lòa. Shino đứng phía sau, giọng cô vẳng lên lần cuối, dịu dàng mà chắc nịch: ``Chúc các cậu may mắn, \emph{Hikawa-san}.''

Cái tên ấy lặp lại, như một dấu ấn khắc sâu, khiến cả ba thoáng khựng lại một nhịp. Có lẽ Sora - hay chính Shino - vẫn còn điều gì muốn nói, nhưng đã để dành cho họ tự mình tìm hiểu.

Ánh sáng trắng nuốt trọn lấy tất cả. Cảnh vật tan biến, âm thanh cũng dần mờ nhạt, chỉ còn lại một sự yên ắng\dots

\dots và một cuộc hành trình đang chờ đợi phía trước ba người họ.

Ánh sáng dần tan, chỉ để lại khoảng trống lặng im trong ngôi đền. Những cánh hoa bỉ ngạn rung rinh theo gió, như thì thầm điều gì đó chưa kịp nói ra.

Ở một nơi khác, một thời khắc khác, những mảnh ký ức vẫn đang bị giằng xé.

Con đường mà Reiko, Saikyou và Kyuen lựa chọn không chỉ dẫn đến sự thật về Sora, mà còn đến cả lời nguyền nơi nhà ga Aogami, nơi đã chôn vùi bao nhiêu bóng hình mất tích.

\end{document}